\newif\ifnotes
\notestrue
%\notesfalse

%\documentclass{article}
\documentclass[runningheads]{llncs}
%\usepackage{geometry}
% \geometry{
%   a4paper,         % or letterpaper
%   textwidth=13cm,  % llncs has 12.2cm
%   textheight=24cm, % llncs has 19.3cm
%   heightrounded,   % integer number of lines
%   hratio=1:1,      % horizontally centered
%   vratio=2:3,      % not vertically centered
% }
\usepackage{graphicx,svg} % Required for inserting images
%\usepackage{background} % for custom headers 
\usepackage{amsmath,amssymb}
\usepackage{mathalpha}
%\usepackage[scr=boondox,  % heavily sloped
%            cal=esstix]   % slightly sloped
%           {mathalpha}
\usepackage{algorithm}
\usepackage[noend]{algpseudocode}
\usepackage[n,advantage,operators,sets,adversary,landau,probability,notions,logic,ff,mm,primitives,events, complexity,asymptotics,keys]{cryptocode}
\usepackage{mathrsfs}
\usepackage{tikz-cd}
\usetikzlibrary{decorations.pathreplacing,angles,quotes}
\usepackage{color}
%\usepackage[table]{xcolor}
\usepackage{colortbl}
\usepackage{url,multirow}
\usepackage{todonotes, xspace}
\usepackage{graphicx}
\graphicspath{{images/}}

\usepackage[bookmarks, 
    colorlinks=false, 
    pdftitle={S-two Whitepaper}, 
    pdfauthor={StarkWare Team}
]{hyperref}

%comment out watermarks when not needed
% \usepackage{draftwatermark}
% \SetWatermarkScale{7}	
% \SetWatermarkColor[gray]{0.9}

% Theorem environments %%%%%%%%%%%%%%%%%%%%%%%%%%%%%%%%%%%%%%%%%%

\spnewtheorem{thm}{Theorem}{\bfseries}{\itshape}
\spnewtheorem*{thm*}{Theorem}{\bfseries}{\itshape}
\spnewtheorem{cor}{Corollary}{\bfseries}{\itshape}
\spnewtheorem{lem}{Lemma}{\bfseries}{\itshape}
\spnewtheorem{prop}{Proposition}{\bfseries}{\itshape}
\spnewtheorem{protocol}{Protocol}{\bfseries}{\itshape}

\spnewtheorem{defn}{Definition}{\bfseries}{}

\spnewtheorem{conj}{Conjecture}{\bfseries}{\itshape}
\spnewtheorem{rem}{Remark}{\bfseries}{}


\newcommand{\m}{\textsf{M31}}
\newcommand{\cm}{\textsf{CM31}}
\renewcommand{\i}{i}
\newcommand{\qm}{\textsf{QM31}}
\renewcommand{\j}{j}
\newcommand{\feltLarge}{\textsf{F252}}

%%%
\usepackage{lipsum}

\def\frontmatter{%
    \pagenumbering{roman}
    \setcounter{page}{1}
    \renewcommand{\thesection}{\Roman{section}}
}%

\def\mainmatter{%
    \pagenumbering{arabic}
    \setcounter{page}{1}
    \setcounter{section}{0}
    \renewcommand{\thesection}{\arabic{section}}
}%

\def\backmatter{%
    \setcounter{section}{0}
    \renewcommand{\thesection}{\Alph{section}}
}%




% \def\companylogo{\includesvg[width=4cm]{StarkWare-logo-01.svg}}
% \fancypagestyle{companypagestyle}
% {
%     \fancyhf{}
%     % \fancyfoot[L]
%     % {
%     %     \parbox[b]{\dimexpr\linewidth-2.5cm\relax}
%     %     {
%     %         \sffamily \bfseries\hfill MyCompany, MyStreet 1\\
%     %         Page \thepage\ of \pageref{LastPage}\hfill 1234 City, Country\\
%     %         {\color{red}\rule{\dimexpr\linewidth\relax}{0.4pt}}\\
%     %         www.mycompany.com
%     %     }
%     % }
%     \fancyfoot[R]
%     {
%         \parbox[b]{2cm}{\companylogo}%
%     }
% }
% \backgroundsetup{
%   scale=1,
%   color=black,
%   opacity=1,
%   angle=0,
%   position=current page.south,
%   vshift=80pt,
%   contents={%
%   \includesvg[width=4cm]{StarkWare-logo-01.svg}
%   }
% }

\title{%
Provekit Whitepaper
}

\authorrunning{World Foundation}

\date{\today}


\setcounter{tocdepth}{2}

% \AddToHookNext{shipout/background}{
%   \begin{tikzpicture}[remember picture,overlay]
%   \node at ([yshift=100pt]current page.south)  {
%     \includegraphics[width=4cm]{images/StarkWare-logo-01.pdf}
%   };
%   \end{tikzpicture}
% }
\newcommand{\RS}{\mathsf{RS}}
\newcommand{\msk}{\mathsf{msk}}
\newcommand{\comm}{\mathsf{Comm}}
\newcommand{\fold}{\mathsf{fold}}
\newcommand{\h}{\mathsf{h}}

\newcommand{\ulrich}[1]{\ifnotes{\small\color{blue} Ulrich: #1}\xspace\fi}
\newcommand{\remco}[1]{\ifnotes{\small\color{teal} Remco: #1}\xspace\fi}
\newcommand{\marcin}[1]{\ifnotes{\small\color{orange} Marcin: #1}\xspace\fi}
\newcommand{\albert}[1]{\ifnotes{\small\color{cyan} Albert: #1}\xspace\fi}

%%%
\usepackage{lipsum}

\title{Provekit R1CS prover}
% \author{Ulrich Hab{\"o}ck\textsuperscript{1}
% %\\
% %\texttt{ulrich.haboeck@gmail.com},
% }
\date{\today}


% macros

\makeatletter
% \newcommand\frontmatter{%
%    \cleardoublepage
%   \pagenumbering{roman}
% }
% \newcommand\mainmatter{%
%   \pagenumbering{arabic}
% }
% \newcommand\backmatter{%
%    \cleardoublepage
%    %\pagenumbering{Roman}
% }
\makeatother



\begin{document}
\let\oldaddcontentsline\addcontentsline
\def\addcontentsline#1#2#3{}
\maketitle
\def\addcontentsline#1#2#3{\oldaddcontentsline{#1}{#2}{#3}}

\frontmatter



% \def\addcontentsline#1#2#3{}
% \section*{Intro}
% \def\addcontentsline#1#2#3{\oldaddcontentsline{#1}{#2}{#3}}
% %\NoBgThispage
% \pagestyle{plain}



%\section{}

\tableofcontents

\mainmatter


\section{R1CS with scheduled witnesses. \textcolor{red}{(tbc)}}
Two-round R1CS, to support lookups into tables, which are baked into the R1CS matrices.
The witness $\vec w$ consists of 
\begin{enumerate}
    \item 
    Round-1 witness $\vec w_1$: 
    These witnesses are independent of the verifier challenge
    Includes public inputs, constants, and lookup table multiplicities.
    \item 
    Round-2 witness $\vec w_2$: 
    These witnesses depend on the verifier challenges.
    These are the verifier challenges themselves, the logUp inverses and their sums.
\end{enumerate}
\ulrich{Let us elaborate on the usefulness of 2-stage R1CS beyond lookups.}
\ulrich{%
Need to define $\mathscr R_{main}$ and $\mathscr R_{aux}$.
Target $\vec w_1\in \mathscr R_{main} \cap \mathscr R_{aux}$
}
% The verifier challenges for logUp are dedicated inputs of the circuit, and the matrices $A$, $B$, $C$ cover the logUp logic. 
The complete witness for the constraint system is of the form 
\begin{equation*}
%\label{e:witness:split}
\vec w = \vec w_1\| \vec w_2,
\end{equation*}
with $\vec w_1\in \mathbb F_q^{N_1}$ and $\vec w_2\in \mathbb F_q^{N_2}$, 
\ulrich{Concatenation is never used anywhere.}
\ulrich{%
It seems more efficient to treat logup with several basefield challenges.
This is for sure the case for the witness builder.
Memory constraints in the zerocheck (due to the higher-dimensional hypercube) can be overcome by univariate skip.
}
subject to $M\geq 1$ constraints
\begin{equation*}
% \label{e:R1CS}
(A\cdot \vec w) \odot (B\cdot \vec w) = C\cdot \vec w,
\end{equation*}
with $A, B, C\in \mathbb F^{K\times N}$, and where $\odot$ is the entrywise (Hadamard) product.
As for the witness, we split the matrices $M= A, B, C$ into 
\begin{equation*}
M = 
\begin{pmatrix}
M_1 | M_2
\end{pmatrix}, 
\end{equation*}
with $M_1\in \mathbb F_q^{K \times N_1}$ and $M_2\in \mathbb F_q^{K\times N_2}$ corresponding to the $\vec w_1$- and $\vec w_2$-parts. 
% \ulrich{Would $M_2$ over $F$ make sense?}
Typically the system is partially decoupled, with each $A_2, B_2, C_2$ starting with a large block of zero rows, but we currently do not leverage simpler representations than the ones above. 
% \ulrich{Do we?}

% \[
% \left[
% \left(
% \begin{array}{c|c}
% A_{1,1}  & 0
% \\\hline
% A_{2,1} & A_{2,2} 
% \end{array}
% \right) 
% \cdot
% \begin{pmatrix}
% \vec w_1 \\ \vec w_2
% \end{pmatrix}
% \right]
% \odot
% \left[
% \left(
% \begin{array}{c|c}
% B_{1,1}  & 0
% \\\hline
% B_{2,1} & B_{2,2} 
% \end{array}
% \right) 
% \cdot
% \begin{pmatrix}
% \vec w_1 \\ \vec w_2
% \end{pmatrix}
% \right]
% =
% \left(
% \begin{array}{c|c}
% C_{1,1}  & 0
% \\\hline
% C_{2,1} & C_{2,2} 
% \end{array}
% \right) 
% \cdot
% \begin{pmatrix}
% \vec w_1 \\ \vec w_2
% \end{pmatrix}
% \]
% \[
% M = 
% \left(
% \begin{array}{c|c}
% M_{1,1}  & 0
% \\\hline
% M_{2,1} & M_{2,2} 
% \end{array}
% \right) 
% \]
% with $M_{2,1}$ and $M_{2,2}$ very sparse (most of the rows have at most 2 entries).

We need:
\begin{itemize}
    \item 
    Definition of the first stage R1CS with auxiliary lookup relation $\mathscr R_\text{R1CS} \cap \mathscr R_\text{aux}$
    \item
    Definition of randomized R1CS (with two stages), and its soundness error
\end{itemize}

\subsubsection{Split-witness builder.}
There is a split-witness builder for stepwise computation of the complete witness.



\section{Preliminaries}

Although the entire construction described in the writeup applies to arbitrary prime fields, we confine ourselve to the the following choices.
%
Let $\mathbb F_q$ be the Goldilocks prime field \cite{}, with the Solinas prime
\[
q = 2^{64} - 2^{32} + 1
\]
as modulus. 
We denote by $F\geq \mathbb F_q$ any cryptographically large extension, in our case the degree 3 extension $F = \mathbb F_{q^5}$, using...
\ulrich{%
Check which irreducible polynomial is best
}


\subsection{Notation}
Boolean Hypercube $H_n = \{0,1\}^n\subset \mathbb F_q^n$, indexing in natural order via $[0,2^n)$, 
multivariate and multilinear polynomials,
Lagrange polynomial
\begin{equation}
\label{e:eq}
eq(\vec y, \vec\rho) = \prod_{i=1}^n (y_i \cdot \rho_i + (1 - y_i)\cdot (1-\rho_i))
%= \prod_{i=1}^m (1 - y_i - \rho_i + 2 y_i  \rho_i),
\end{equation}

Inner product notation.


\subsection{Multilinears and  univariates}

 
Given a function $f\in \mathbb F^{H_n}$ (a vector of values) we will write
\[
f(X_1,\ldots, X_n) \in \mathbb F[X_1,\ldots, X_n]^{<2}
\]
for its multilinear extension, and 
% multilinear $f(X_1,\ldots, X_n)\in $
% %, given by its values over the Boolean hypercube $H_n$, %(the component of an interleaved representation) with the univariate polynomial
associate to it the univariate polynomial
\begin{equation}
\label{e:values:to:coeffs}
u_f(X) = \sum_{i=0}^{2^n-1} f(i_1,\ldots, i_n) \cdot X^i,
\end{equation}
defined by the values of $f$ over $H_n$,  where $i = i_1 + \ldots + 2^{n-1} i_n$ is the decomposition of the index into bits. 
Under this \textit{values-to-coefficient} regime, we identify the space of multilinears (in $n$ variables) with the univariate space
\begin{equation}
\mathscr P_{n}(\mathbb F) = \mathbb F[X]^{<2^n} 
\end{equation}
in a way that is particularly efficient for Reed--Solomon encoding. 

% \subsubsection{Liftings.}
% Given a multilinear $f$ in $k$ variables, and $n>k$, we call
% \[
% f^*(X_1,\ldots, X_n) = (1-X_1)\cdot\ldots (1-X_{n-k})\cdot f(X_{n-k+1},\ldots, X_n)
% \]
% its \textit{lifting} to $n$ variables. 
% Lifting will allow to treat a collection of polynomials in a uniform manner, while not paying the cost of duplicated values.
% In terms of univariate representations,
% \[
% u_{f^*}(X) = u_f(X^{2^{n-k}}).
% \]

\subsubsection{Univariate evaluation}
Evaluation at a point $z\in\mathbb F$ is a linear functional and can be written as inner product $u_m(z) = \langle f, L_z\rangle_{H_n}$ with  
\begin{equation}
\label{e:uni:eval:kernel}
L_z(X_1, \ldots, X_n) = \prod_{j=1}^n \left(
 (1 - X_j) + X_j \cdot z^{2^{j-1}}  \right),
\end{equation}
and its values over the hypercube are the powers of $z$ in natural order.
% In terms of codes, $L_z$ expresses point evaluation of the codeword as linear functional on the message.

% This establishes a one-to-one correspondence between multilinear polynomials in $k$ variables, and univariate polynomials from the space 
% \[
% \mathscr P_{k} = \mathbb F[X]^{<2^k}.
% \]

% For example the univariate Lagrangian 
% \[
% L_z(X) = \frac{1}{2^n}\cdot \left(1 + \sum_{i=1}^{2^n-1} z^{2^n - i} \cdot X^i\right) 
% \]

% Interleaved representations are encoded component-wise, i.e. row by row, using a Reed--Solomon code with message length matching the row length. 
% Concretely, each postfix function $g(X_1, \ldots, X_{n-k}) = f(X_1,\ldots, X_{n-k},\vec y)$, given by its values over the Boolean hypercube $H_{n-k}$, is associated with the univariate polynomial
% \[
% \sum_{\vec x\in \{0,1\}^{n-k}} g(\vec x) \cdot X^{\iota(\vec x)} \quad \in \mathbb F[X]^{<2^{n-k}},
% \]
% where $\iota(\vec x) = x_1 + 2 x_2 + \ldots + 2^{n-k-1} x_{n-k}$ is the natural index of $\vec x$, and evaluated over a specified domain $D\subset \mathbb F$. 
% In this way, we obtain postfix many Reed--Solomon codes, viewed as word from the \textit{interleaved code}.






\subsection{Interleaved representation}

In alignment with the sumcheck protocol, we will often think of a multilinear $f(X_1,\ldots, X_n)$ decomposed into its $k$-prefix functions
% by postfixes $\vec y\in \{0,1\}^k$,
\begin{align}
\label{e:prefix:decomp}
f(X_1,\ldots,X_n) =  \sum_{i = 0}^{2^{k}-1} eq_{i}(X_{1},\ldots, X_k) \cdot f_i(X_{k+1},\ldots, X_{n}) \cdot 
\end{align}
Here we again used the identification  $H_k = [0,2^k)$, by each index $i$ specifying the prefix $(i_{1},\ldots, i_{k})\in H_{k}$ via its bits.
We call 
$
(f_0, \ldots, f_{2^{n-k}-1}) %\in \mathbb F[X_{k+1},\ldots, X_n]^{2^{k}}
$
the \textit{interleaved representation} of $f$, with interleaving depth $e = 2^{k}$.
In terms of their univariate polynomials,
\begin{align}
\label{e:FFT:decomp}
u_f(X) = \sum_{i=0}^{2^k-1}  X^{i} \cdot u_{f_i}(X^{2^k}), 
\end{align}
which is the FFT decomposition of $u_f(X)$. 

\begin{rem}%[Postfix decomposition]
\label{rem:postfix:decomp}
The current version of ProveKit instead employs a postfix decomposition of multilinear polynomials, which corresponds, in the univariate representation, to partitioning the polynomial into contiguous blocks of monomial powers.
%$u_f(X) = \sum_{i=0}^{2^{n-k}-1} u_{f_i}(X) \cdot X^{i \cdot 2^{k}}$.
Since our construction relies very little on properties specific to Reed--Solomon codes, the choice of decomposition is largely a matter of convenience. We adopt a prefix decomposition instead, as it leads to notation that is more consistent with our treatment of liftings.
\end{rem}



\subsection{Liftings}

Given a multilinear $f$ in $k$ variables, and $n>k$, we call
\begin{equation}
\label{e:lifting:multilinear}
f^*(X_1,\ldots, X_n) = (1 - X_1)\cdot \ldots \cdot (1 - X_{n-k}) \cdot f(X_{n-k+1},\ldots, X_n)
\end{equation}
its \textit{lifting} to $n$ variables. 
In terms of the values over the hypercube, lifting simply places the original vector (the values over $H_k$) on the prefix cube 
\[
\{0^{n-k}\}\times H_{k} \subset H_n,
\]
and pads it with zero on the remainder of $H_n$.
In terms univariate representations, 
\begin{equation}
u_{f^*}(X) = u_f(X^{2^{n-k}}).
\end{equation}
We use lifting to uniformize our R1CS across the two stages at no additional cost.


% \subsection{Sumcheck}

% % In reverse order, since the hypercube is indexed in natural order.

% The multivariate sumcheck protocol \cite{Sumcheck} is an interactive proof for that a multivariate polynomial $G\in \mathbb F[X_1,\ldots, X_n]$ satisfies
% \[
% \sum_{\vec x \in H_n} G(\vec x) = s.
% \]
% % Typically, the multivariate $G$ is a virtual polynomial in several committed multilinears, and its maximum individual degree $d = \max_i \deg_{X_i}(G)$ is at least $2$.
% We assume that the reader is familiar with it, and only point out the following:
% For implementation reasons, we consider the sumcheck protocol in \textit{reverse order}, from the last variable $X_n$ down to the first variable $X_1$.
% In each step,  $k = n, n-1, \ldots, 1$, the prover provides the refinement polynomials $q_k(X)\in \mathbb F[X]^{\leq d_k}$, where $d_k$ is the individual degree of $G$ in $X_k$, and the verifier chooses to continue on the specialized claim $q_k(\lambda_k)$ at a random  $\lambda_k\sample\mathbb F$. 
% \[
% q_k(X) = \sum_{(x_{1},\ldots, x_{k-1}) \in H_{k-1}} G(x_1,\ldots, x_{k-1},X,\lambda_{k+1},\ldots, \lambda_n),
% \]


% Although we shall describe our protocols in a self-contained manner, we quickly sketch the sumcheck protocol, as it is a core ingredient of multilinear proof system.

% Before the first round the prover claims the refinement polynomial 
% \[
% q_n(X) = \sum_{(x_2,\ldots, x_{n-1}) \in H_{n-1}} G(x_1,\ldots, x_{n-1}, X),
% \]
% which is of degree at most $d$.
% The verifier checks the sanity condition that $s = q_0(0) + q_0(1)$, chooses a random $\lambda_n$, and it asks the prover to prove the specialized claim
% \[
% q_n(\lambda_n) = \sum_{(x_2,\ldots, x_n) \in H_{n-1}} G(x_1,\ldots, x_{n-1}, \lambda_n)
% \]
% instead.
% The protocol continues in the same manner, letting the prover provide refinement polynomials
% \[
% q_k(X) = \sum_{(x_{1},\ldots, x_{k-1}) \in H_{k-1}} G(x_1,\ldots, x_{k-1},X,\lambda_{k+1},\ldots, \lambda_n),
% \]
% for $i=1,\ldots,n-2$ as response of the received challenges $\lambda_2,\ldots, \lambda_{n-1}$.
% In the last step the verifier samples $\lambda_{n}\sample F$, which eventually reduces the sumcheck claim to that
% \[
% q_{n-1}(\lambda_n) = G(\lambda_1,\ldots, \lambda_n), 
% \]
% a evaluation claim for the multivariate polynomial $G$ at the random point $\vec \lambda = (\lambda_1,\ldots, \lambda_n)$, comprised of the random challenges received in the course of the protocol.
% For a detailed treatment of the sumcheck protocol, and its importance to multivariate proofs in general, we refer to the manuscript \cite{Thaler:book} and the references therein.



\subsection{Properties of interleaved Reed--Solomon codes}

Are RS codes over an extension field. 
Same rate, same distance.  




\section{The protocol without zero-knowledge}
\label{s:protocol:no:zk}

In this section we describe the main construction for proving a two-stage R1CS, yet without the additional measures for zero-knowledge.
The protocol is a combination Spartan \cite{Spartan} and WHIR \cite{WHIR} as a general inner-product argument. 
We skip the proving the values of the public circuit polynomials, which in \cite{Spartan} is handled by SPARK, and assume that the verifier does compute their values by itself.
In this way we obtain succinct proof sizes, yet the verifier costs are of the same order as the witness size.

%is being outsourced to the verifier, yielding only succinct proof sizes, yet not a succinct verifier.

% In the following we consider the R1CS stage 1 and 2 witnesses $\vec w_1$ and $\vec w_2$ padded to their next power of two, and hence of twoadic size $2^{m_1}$ and $2^{m_2}$.
% Without loss in generality we assume that $m_1 \geq m_2$, and we will view $w_2$ padded to the same length via lifting their MLEs.
% that $2^{m_1}$ matches $2^{k_0}$, the dimension of the largest function space $\mathscr P_{k_0}$ as specified by WHIR.


\subsection{Arithmetization.}
% The R1CS stage 1 and 2 witness vectors $\vec w_1, \vec w_2$ are zero-padded to their next power of two, and placed on the Boolean hypercube of respective size in the common manner.  
We consider the R1CS witness vectors padded to their next power of two, i.e. as functions
\[
w_i: \{0,1\}^{m_i} \longrightarrow \mathbb F_q, \quad i=1,2,
\]
%and $w: \{0,1\}^n \longrightarrow F$, 
and we denote their multilinear extensions again by the same letters.
Likewise, the number of overall constraints are padded to a power of two, say $2^m$, and we view the R1CS matrices as functions
\[
A_i, B_i, C_i: \{0,1\}^m \times \{0,1\}^{m_i} \longrightarrow \mathbb F_q, \quad i =1,2, 
\]
with the same notation for their multilinear extensions. 

Throughout the following we assume that $m_1\geq m_2$. 
% \subsubsection{Lifted R1CS.}
We lift $w_2$ to the same number of variables as $w_1$, and likewise we do for the second stage matrix polynomials $A_2, B_2, C_2$.
The resulting system has
\begin{multline}
\label{e:combined:witness:poly}
w(X_0,X_1,\ldots, X_{m_1}) =
\\
(1- X_0) \cdot w_1(X_1, \ldots, X_{m_1}) + X_0 \cdot w_2^*(X_{1},\ldots, X_{m_1}),
\end{multline}
as the combined witness polynomial, 
and
\begin{multline}
\label{e:combined:matrix:poly}
M(\vec Y,X_0,X_1,\ldots, X_{m_1}) =
\\
(1- X_0) \cdot M_1(\vec Y, X_1, \ldots, X_{m_1}) +
 X_0 \cdot  M_2^*(\vec Y, X_1,\ldots, X_{m_1}),
\end{multline}
for $M=A,B,C$, as the combined matrix polynomials.
%where $\vec Y=(Y_1,\ldots, Y_m)$, and $M_2^*(\vec Y,\:.\:)$ is the lifting of $M_2(\vec Y,\:.\:)$ to $m_1$ variables.


\ulrich{placing of inputs}
% Again, the multilinear extensions of the complete matrices $M$ and their components $M_1, M_2$ are connected via
% \begin{multline*}
% M(\vec Y, X_1,\ldots,X_n) = (1 - X_n)\cdot M_1(\vec Y, X_1, \ldots, X_{n-1}) 
% \\ + X_n \cdot M_2(\vec Y, X_1, \ldots, X_{n-1}).
% \end{multline*}



\subsection{Encoding.}

We encode $w_1$ and $w_2^*$ in interleaved manner, using the same basecode
\begin{equation}
\label{e:base:code}
C = \RS[F,D, 2^{k}],
\end{equation}
which corresponds to the first basecode of WHIR.
(We will discuss the setup of WHIR later.)
The multilinears $w_1$ and $w_2^*$ are decomposed into $(m_1 - k)$-prefix components, 
%in the $k$ variables $X_{m_1-k+1}, \ldots$, $X_{m_1}$, 
% \[
% w_{i,0}(X_{m_1-k+1},\ldots, X_{m_1}),\ldots, w_{i,e_i-1}(X_{m_1-k+1},\ldots, X_{m_1})
% \]
%as in \eqref{e:prefix:decomp}, 
which are then encoded non-systematically, using their univariate representation from $\mathscr P_{k}$.
%in the previously described values-to-coefficient regime.
% \footnotemark
% \footnotetext{%
% Given the values $(w_1(i))_{i\in H_{m_1}}$, we reshape it into a matrix of $2^{m_1-k_1}$ rows of length $2^{k_1}$, $(w_1(i + 2^{m_1-k_1}j))_{i,j}$,
% and encode each of the rows separately, taking the row entries as coefficients of the polynomials.
% }
The resulting words 
\[
f_1, f_2^* \in C^{e},
\]
the interleaved Reed--Solomon code of interleaving depth $e = 2^{m_1-k_1}$, 
where $f_2^*$ has only $d = 2^{\max\{m_2-k_1,0\}}$ non-zero components. 
(These are the prefix embeddings of the $d$ components of $w_2$.)
Thus committing to the lifted $w_2^*$ causes no overhead compared to $w_2$, given that $m_2 \geq k$.
% Whenever $m_2$ is smaller than $k$, we will work with lifted commitments.


% assuming that $m_i\geq k$. 


% This amounts to decomposing the combined witness polynomial into 
% \begin{align*}
% w(X_1,\ldots, X_{m_1 + 1}) = w(X_1,\ldots, X_k, \vec y)    
% \end{align*}



% While the setup of WHIR will be so that $m_1 \geq k$, which is typically the larger witness, it may happen that the second stage witness is smaller than the message length of $C$. 
% In this case we use \textit{lifting}, which we decribe next.

% \ulrich{%
% More elegant: Let us lift $p_2(X)$ to $p_2(X^{e})\in\mathscr P_{m_1}$ with $e=2^{m_1 - m_2}$, and decompose it into $\mathscr P_{k_1}$-components.
% If $m_2 < k_1$, i.e. $e > e_1$, we should end up with the lifting as described below:
% There $p_2^*(X) = p_2(X^{e_1 \cdot 2^{k_2 - m_2}}) $
% \[
% p_2(X^e) = p_2((X^{e_1})^{e/e_1})
% \]
% gives
% \[
% p_{2,1}(Y) = p_2(Y^{e/e_1}) = p_2(Y^{2^{m_1-m_2-m1+k_1}}) = p_2(Y^{2^{k_1 - m_2}}).
% \]
% But we do not want to commit that over $D_1$.
% }


\subsubsection{Lifting small commitments.}
Whenever $m_2 < k$, we commit the univariate representation $u_{2}(X)$ of $w_2$ over a projection of $D$, the image of the multiplicative homomorphism
\begin{equation}
\label{e:projection}
\pi: D \longrightarrow \pi(D),\quad x\mapsto x^{2^{k_1 - m_2}}.
\end{equation}
Since $D$ is an FFT domain, the map is $2^{k - m_2}$-to-$1$ and onto.
The committed word $f_2 \in \mathbb F^{D'}$ is from the smaller code
\begin{equation}
\label{e:projected:code}
C' := \RS[F, D', 2^{m_2}],
\end{equation}
which has the same relative rate as $C$.
The lifted word
\begin{equation}
\label{e:lifting:uni}
f_2^*(x) =  f_2(\pi(x)) %= f_2 (x^{2^{k-m_2}})
\end{equation}
belongs to $C$, since $u_2^*(X) = u_2(X^{2^{k-m_2}})$ is of degree less than $2^{k}$.
%
By means of the projection $\pi$ we simulate $f_2^*\in C$ without paying the price of committing over the larger domain.


\begin{rem}
As noted in Remark~\ref{rem:postfix:decomp}, the current implementation employs a postfix decomposition. However, liftings of the corresponding univariate polynomials are naturally expressed in terms of prefix embeddings, resulting in a somewhat inelegant mixture of prefix and postfix terminology.
\end{rem}

% \ulrich{Ref to Ligerito and zook}
%(See Section \ref{s:WHIR} for details.)

% We use the same commit principle on application layer (i.e. on top of WHIR) for polynomials from \textit{any} $\mathscr P = \mathscr P_k$, 
% \[
% \mathscr P_0 \supseteq \mathscr P \supsetneq \mathscr P_{k_r},
% \]
% not necessarily matching one of the spaces in \eqref{e:chain:function:spaces}: 
% %as interleaved codeword with one of the base codes $C_1,\ldots, C_r$. 
% %from one of the spaces in \eqref{e:chain:function:spaces}.
% Take the unique $i$ so that $\mathscr P_{k_i}\supseteq \mathscr P \supsetneq \mathscr P_{k_{i+1}}$,
% and denote the co-dimension of $\mathscr P_{k_{i+1}}$ in $\mathscr P$ by
% \begin{equation}
% e = \dim(\mathscr P/\mathscr P_{k_{i+1}}).
% %= 2^{k_{i} - k_{i+1}},
% \end{equation}
% The FFT-decomposition of any $p(X) \in \mathscr P$ into 
% %a vector of $e_i$ polynomials $p_1(X),\ldots, p_{e_i}(X)$ from $\mathscr P_{i+1}$,
% \begin{align}
% \label{e:FFT:decomp}
% p(X) = p_1(X^e) + X\cdot p_2(X^e) + \ldots + X^{e - 1}\cdot p_{e}(X^e),
% \end{align}
% defines a linear identity
% \begin{equation}
% \label{e:FFT:decomp:identity}
% \mathscr P \simeq \mathscr P_{k_{i+1}}^{e},    
% \end{equation}
% %
% and with this identity we commit $p(X)\in\mathscr P$ as word from 
% \begin{align}
% %\label{e:C_i}
% C_{k_{i+1}}^e 
% &= 
% \left\{
% [c_1(x),\ldots, c_{e}(x)]\big|_{x\in D_{i+1}} \::\: \text{all }c_j \in C_{k_{i+1}}
% \right\}.
% \end{align}
%\ulrich{Maybe}
% The codes $C_1,\ldots, C_r$ generated by the spaces $\mathscr P_{k_1}, \ldots \mathscr P_{k_r}$ over the domains $D_1, \ldots, D_r$, are the base codes of our scheme. 
% Their minimum distance is
% \begin{equation}
% \label{e:min:dist:baseCode}
% \delta(C_i) = 1 - \rho(C_i) - \frac{1}{|D_i|},  
% \end{equation}
% where
% \begin{equation}
% \label{e:rate:baseCode}
% \rho(C_i) = \frac{2^{k_{i}}}{|D_i|} = 2^{k_{i}- n_{i}}
% \end{equation}
% is the relative rate.
% \ulrich{%
% Maybe add blow-up factors?
% }





\subsection{The IOPP}


%Both witnesses will be committed as described above, as interleaved codeword with base code one of the $C_1,\ldots, C_r$ from \eqref{e:base:codes}.

% In many of our applications, the second stage witness is not much smaller than the first, $m_1 = m_2 + 1$, but we will not make use of this fact.


%The complete witness $\vec w = \vec w_1 \| \vec w_2$ is then of size $N=2^n$.


% The WHIR function spaces $\mathscr P_{k_0}$, \ldots $\mathscr P_{k_r}$ can be chosen freely, except that $k_0 = m_1$, so that $\mathscr P_{k_0}$ matches the size of the witness $w_1$.
% The size of $w_2$ either matches the dimension of one of these spaces, or an intermediate power of them. 
% Take $j_2 \in \{k_1, \ldots, k_r\} $ as the unique value so that $m_2 \in (k_{j_2}, k_{j_2 - 1}]$. 
% Then the univariate representation
% \[
% w_2(X) \in \mathscr P_{k_{j_2}}^{e_{j_2}'},
% \]
% with intermediate exponent $e_{j_2}' = 2^{m_2 - k_{j_2}}$ from $(1, e_{j_2}]$. 
% % 2^m_2 = e_{j_2}' \cdot 2^{k_{j_2}}
% We commit it as word from the interleaved code $C_{k_{j_2}}^{e_{j_2}'}$. 
% \ulrich{This needs to be moved to function spaces and codes.}





% \subsubsection{Commitments.}
% First stage witness $w_1$ is committed as codeword from $C_1^{e_0}$, i.e. $k_1\cdot e_1 = 2^{n_1}$ and 
% $w_2$ is committed as codeword from one of the $C_i$, with interleaving depth as in $C_i$, or less (as long as a power of two).



% The witnesses $w_1$ and $w_2$ are committed as Reed-Solomon codewords from
% \[
% C_{i} = \RS[F, D_i, 2^{n_i}] = \{p(x)|_{x\in D_i} \::\: p \in F[X]^{<2^{n_i}} \}, \quad i=1,2,
% \]
% respectively, of possible different rate, and so that their evaluation domains are WHIR compatible, i.e. $D_2 \cap \pi^{n_1-n_2}(D_1) =\varnothing$ whenever $n_2 < n_1$.
% \ulrich{%
% Not sure if this is the best setting. 
% This depends very much how we take the different withness lengths into account in WHIR. 
% I think that a constant rate, and projection-consistent domain, choice is the better way to present it here. 
% }




% \subsubsection{The proof of proximity.}
Given the commitment of $w_1$ as interleaved word $f_1$ on $D$, and proximity parameter $\theta\in (0,1)$, the target of the protocol is to show that $f_1$ is $\theta$-close to a solution of our augmented constraint system, characterized by  $\mathscr R_{main}$ and $\mathscr R_{aux}$ and the inputs $\vec{in} = (in_1,\ldots, in_{t_1})$.
%
% Recall that we use the values-to-coefficients regime, hence 
That is, we want to prove that the committed function belongs to
\begin{equation}
\label{e:IOPP:main:R:0}
\mathcal R_\theta = \left\{
    f\in \left(\mathbb F_q^{D}\right)^{e} \::\:
    \begin{matrix}
    \exists p(X) = \sum_{i=0}^{2^{m_1}- 1} c_i X^i,  %\in \mathbb F_q[X]^{<2^{m_1}}, 
    d(p, f) \leq \theta,
    \\
    \text{ with }
    \\
    (c_i)  \in \mathscr R_{main} \cap \mathscr R_{aux},
    \\
    (c_1,\ldots,c_{t_1}) = \vec{in}
    \end{matrix}
\right\},
\end{equation}
where $d(p,f)$ refers to the interleaved representation of $p(X)$ as  $e = 2^{m_1-k}$ component polynomials of degree $< 2^{k}$.
The protocol is divided into three phases.

\paragraph{Phase 1.}
The first phase is devoted to the second stage of the R1CS.
The prover receives the stage-2 challenge $\vec\xi\sample\mathbb F_q^{t_2}$ from the verifier, and responds with the commitment of the lifted witness $w_2^*(X_1,\ldots, X_{m_1})$.
% In this phase we apply the well-known DEEP principle \cite{DEEPFRI} to simplify the soundness analysis in the heterogeneous setting, where the base codes for $\vec w_1$ and $\vec w_2$ are different.  
% We ask the prover for the value of the committed polynomials at random out-of-domain points $z_1$ and $z_2$, respectively.
% Then with high probability, the claimed values narrow down the prover to a single proximate polynomial for each of the commitments, mitigating quadratic expressions in the list size. 
% (See below for a more in-depth discussion of the technique.)
% %Since we only deal with two witnesses, the impact on proof sizes is negligible.
% %
% The correctness of the claims, rephrased as the inner-product constraints 
% %for the univariate opening claims
% \begin{equation}
% \label{e:witnesses:DEEP:eval}
% v_i = \sum_{\vec x \in H_{m_i}} L_{z_i}(\vec x)\cdot w_i(\vec x), \quad i = 1,2,
% \end{equation}
% with the multivariate representation of the univariate Lagrangians (recall \eqref{e:uni:Lagrangian:multi:rep}) are attested later on by  WHIR.
% (Here $L_{z_i}$ is the multilinear representation of the univariate Lagrangian ź, see equation \eqref{e:uni:Lagrangian:multi:rep}.
%(We use the same notation for the two.)

% \begin{equation}
% \label{e:IOPP:main:R:phase:1}
% \mathcal R_{1} = \left\{
%     f\in \left(\mathbb F_q^{D_1}\right)^{e_1} \::\:
%     \begin{matrix}
%     \exists p(X) = \sum_{i=0}^{2^{m_1}-1} c_i \cdot X^i \in \mathscr P_{m_1}, 
%     \\
%     d((p_1,\ldots, p_{e_1}), (f_1,\ldots, f_{e_1})) \leq \theta,
%     \\
%     % \text{coeffs of }p(X) 
%     (c_i)\in \mathscr R_{R1CS} \cap \mathscr R_{lookup}
%     \end{matrix}
% \right\},
% \end{equation}

\paragraph{Phase 2.}
The Spartan zero-check.
We prove constraint satisfaction of the multilinears
\begin{align}
\label{e:a:poly}
w_M(\vec Y) &= 
\sum_{\vec x\in H_{m_1+1}} M(\vec Y, \vec x)\cdot w(\vec x), 
\\\notag
&= \sum_{\vec x\in H_{m_1}} M_1(\vec Y, \vec x)\cdot w_1(\vec x) + \sum_{\vec x\in H_{m_1}} M_2^*(\vec Y,\vec x) \cdot w_2^*(\vec x)
\\\notag
&= \sum_{\vec x\in H_{m_2}} M_1(\vec Y, \vec x)\cdot w_1(\vec x) + \sum_{\vec x\in H_{m_2}} M_2(\vec Y,\vec x) \cdot w_2(\vec x),
% \quad 
% b(\vec Y) = 
% \sum_{\vec x\in H_{m_1+ 1}} B(\vec Y, \vec x)\cdot w(\vec x), 
% \\
% c(\vec Y) =
% \sum_{\vec x\in H_{m_1+1}} C(\vec Y, \vec x)\cdot w(\vec x),
\end{align}
for $M=A,B,C$.
% \begin{equation}
% \label{e:a:poly}
% \begin{aligned}
% a(\vec Y) = 
% \sum_{\vec x\in H_{m_1+1}} A(\vec Y, \vec x)\cdot w(\vec x), 
% \quad 
% b(\vec Y) = 
% \sum_{\vec x\in H_{m_1+ 1}} B(\vec Y, \vec x)\cdot w(\vec x), 
% \\
% c(\vec Y) =
% \sum_{\vec x\in H_{m_1+1}} C(\vec Y, \vec x)\cdot w(\vec x),
% \end{aligned}
% \end{equation}
%which are uniquely determined by the commitments of $w_1$ and $w_2^*$.
%
For showing that
\begin{equation*}
% \label{e:R1CS:domain:identity}
w_A(\vec y)\cdot w_B(\vec y) - w_C(\vec y) = 0, \quad \text{for all }\vec y \in H_m.
\end{equation*}
we run the sumcheck protocol on its inner product with $eq(\:.\:,\vec\rho)$, where $\vec\rho\sample F^m$ is random.
% %of the quadratic expression $a(\vec y)\cdot b(\vec y) - c(\vec y)$ with 
% \begin{equation*}
% % \label{e:R1CS:zerocheck}
% \sum_{\vec y\in H_m} \left(a(\vec y)\cdot b(\vec y) - c(\vec y)\right) \cdot eq(\vec y, \vec \rho) = 0,
% \end{equation*}
% where $eq(\vec y,\vec\rho)$ is 
% the $m$-dimensional Lagrangian $eq(\vec y, \vec\rho)$ at a random $\vec \rho\sample F^{m}$ sampled by the verifier.
% The sumcheck protocol 
This reduces the above claim to evaluation claims 
\begin{equation}
\label{e:image:polys:evals}
v_a = w_A(\vec \alpha),\quad v_b = w_B(\vec\alpha),\quad v_c = w_C(\vec\alpha),
\end{equation}
where  $\vec\alpha\in F^m$ are the verifier challenges drawn in the course of the sumcheck rounds.
% (The value $eq(\vec\alpha, \vec \rho)$ is public and can be computed cheaply.)
%and efficiently computable in $O(m)$ steps.) 

% To prove the claims
% \begin{align*}
% v_a &= \sum_{\vec x\in H_n} A(\vec \alpha, \vec x)\cdot w(\vec x), \quad v_b = \sum_{\vec x\in H_n} B(\vec \alpha, \vec x)\cdot w(\vec x)
% \\
% v_c &= \sum_{\vec x\in H_n} C(\vec \alpha, \vec x)\cdot w(\vec x),
% \end{align*}
% together with the consistency constraints with respect to the public circuit inputs, including constants and the logUp challenges,
% \begin{align*}
% \vec {in}_j &= \sum_{\vec x\in H_n} eq(j, \vec x)\cdot w(\vec x), \quad j=0,\ldots,I,
% \\
% \alpha_j &= \sum_{\vec x\in H_n} eq(2^{n-1} + j,\vec x) \cdot w(\vec x), \quad j= 0,\ldots,k
% \end{align*}

\paragraph{Phase 3.}
The third phase of the protocol is devoted to proving the linear claims \eqref{e:image:polys:evals}, which is done by WHIR on the combined witness $w$.
For public input and challenge consistency, we further add the constraints
\begin{align}
\label{e:input:w1:constraint}
in_j &= \sum_{\vec x\in H_{m_1}} eq(j,\vec x)\cdot w_1(\vec x), \quad j=0,\ldots, t_1,
\\
\label{e:input:w2:constraint}
\xi_j & = \sum_{\vec x\in H_{m_1}} eq^*(j,\vec x)\cdot w_2^*(\vec x), \quad j=0,\ldots, t_2,
\end{align}
% \albert{Where are these alpha's in the prefix of $w_2$ coming from? Oh ok they are logUp challenges. Why is there a tuple of challenges and not one?} 
% \ulrich{%
% Usually, you have one logUp challenge for the point, where the fractions are evaluated at, and other challenges for compressing vectors into scalars, and sometimes even for domain separation (but I do not know whether we do the latter with randomness).
% }
where $\vec\xi = (\xi_1,\ldots, \xi_{t_2})$ are the challenges drawn in the second stage of the R1CS.
The WHIR proof of proximity eventually reduces the linear claims on $w$ to claims on the combined matrix polynomials, which need to be checked by the verifier at full expense.

Let us give the formal description of our IOPP. 
For simplicity we assume that 
\[
t_1 = t_2 = t.
\]
Recall that $m_1\geq m_2$.

% \begin{equation}
% \label{e:split:matrix:values}
% A(\vec\alpha,\vec\beta), 
% \quad 
% B(\vec\alpha,\vec\beta), 
% \quad 
% C(\vec\alpha,\vec\beta),
% \end{equation}
%and those of the involved Lagrangians $L_z(\vec\beta)$ and $eq(j,\vec \beta)$, 


% As we emphasized before, we do not target a succinct verifier and leave the verification of \eqref{e:split:matrix:values} to the verifier, instead of providing a separate proof for it.


\begin{protocol}[Main IOPP]
\label{prot:main}
The R1CS matrices $A_i,B_i,C_i \in \mathbb F_q^{m\times m_i}$, $i=1,2$, and public inputs $\vec{in} = (in_0,\ldots, in_{t})\in \mathbb F_q^t$ are known to both the prover and the verifier. 
The prover has committed to the first-stage witness $\vec w_1\in \mathbb F_q^{m_1}$ as interleaved Reed--Solomon codeword $f_1\in C^e$ as described above, and provided the verifier oracle access.
\begin{enumerate}
\item
\label{i:main:phase:1}
R1CS second stage:
% \begin{enumerate}
% \item
% \label{i:main:DEEP:w1}
% The verifier sends a random $z_1\sample F$ to the prover, which answers with the value $v_1$ of the univariate representation $w_1(X)$ at $z_1$.
% %
% \item
% \label{i:main:logUp}
The verifier samples the stage-2 challenge
\[
\vec\xi \sample \mathbb F_q^{t}
\]
and sends them to the prover.
The prover computes $\vec w_2 \in \mathbb F_q^{m_2}$  the second-stage witness of the R1CS, commits to it as codeword $f_2\in C^e$, as described above, and provides the verifier oracle access.
%(the logUp inverses and their sums), 
% commits to it as interleaved Reed--Solomon codeword $f_2$ associated with $\mathscr P_{m_2}(\mathbb F_q)$, and provides the verifier oracle access. 
%
% \item 
% \label{i:main:DEEP:w2}
% The verifier sends a random $z_2\sample F$ to the prover, which answers with the value of the univariate representation $w_2(X)$ at $z_2$.
% \end{enumerate}
%answers with the logUp R1CS witnesses $\vec w_2$, so that the 

\item 
\label{i:main:phase:2}
R1CS zero-check:
% Both prover and verifier engage in a \textit{zero-check} for the domain identity
% \begin{equation}
% a(\vec x) \cdot b(\vec x) = c(\vec x), \quad \vec x\in H_n = \{0,1\}^n
% \end{equation}
% where $a, b$ and $c$ are the multilinear extensions of the implicitly committed matrix-images
% \[
% \vec a = A\cdot \vec w, \vec b = B\cdot \vec w, \vec c = C\cdot \vec w.
% \]
The verifier sends a random 
\[
\vec\rho \sample F^m,
\]
and both prover and verifier engage in the sumcheck protocol for %the R1CS constraints
\begin{align}
\label{e:R1CS:zerocheck}
\sum_{\vec x\in H_m} \left(w_A(\vec x)\cdot w_B(\vec x) - w_C(\vec x)\right)\cdot eq(\vec x, \vec\rho) = 0,
\end{align}
where $w_A(\vec x), w_B(\vec x), w_C(\vec x)$ are the multilinears as specified in \eqref{e:a:poly}.
After the last round of the sumcheck, the prover sends 
% (The protocol reduces the claim on the hypercube sum \eqref{e:R1CS:zerocheck} to a quadratic check on the evaluation claims 
\begin{equation*}
v_A = w_A(\vec\alpha),  \quad v_B = w_B(\vec\alpha), \quad v_C = w_C(\vec\alpha),
\end{equation*}
where $\vec\alpha \in F^m$ are the verifier challenges drawn in the course of the sumcheck, and the verifier checks if 
$q_m(\alpha_m) = (v_A\cdot v_B - v_C) \cdot eq(\alpha,\rho)$.
%These claims are announced by the prover in the last round. 

\item 
\label{i:main:phase:3}
WHIR: 
% The verifier samples a random $\lambda\sample F$, 
Both prover and verifier run WHIR for the following linear constraints on the combined witness polynomial $w$:
\begin{align}
\label{e:phase3:claim:vM}
v_M = \sum_{\vec x\in H_{m_1}} M_1(\vec\alpha, \vec x) \cdot w_1(\vec x) + \sum_{\vec x\in H_{m_1}} M_2^*(\vec\alpha,\vec x)\cdot w_2^*(\vec x)
\end{align}
with $M = A,B,C$, and
\begin{align}
\label{e:phase3:claim:inputs}
in_i = \sum_{\vec x \in H_{m_1}} eq(i, \vec x) \cdot w_1(\vec x),
\quad 
\xi_i = \sum_{\vec x \in H_{m_2}} eq^*(i, \vec x) \cdot w_2^*(\vec x),
\end{align}
for $i = 0,\ldots, t$.
% (The concrete form of the linear combination is described below.)
(This phase reduces the validity of the linear constraints to claims on the combined matrix polynomials
\begin{equation}
\label{e:split:matrix:values}
A(\vec\alpha,\vec\beta), 
\quad 
B(\vec\alpha,\vec\beta), 
\quad 
C(\vec\alpha,\vec\beta),
\end{equation}
where $\vec\beta \in F^{m_1}$ are the folding randomnesses drawn in the course of WHIR.)
\end{enumerate}
If all verifier checks in Phase \ref{i:main:phase:2} and Phase \ref{i:main:phase:3}, including that of the final matrix values \eqref{e:split:matrix:values}, pass, then the verifier accepts.
Otherwise, it rejects.
\end{protocol}

% The three claims \eqref{e:phase3:claim:vM} and \eqref{e:phase3:claim:inputs} are of inner products of $w$ 
% with the kernels
% \[
% K_M = (1-X_0) \cdot M_1(\vec\alpha,X_1,\ldots, X_{m_1}) + X_0\cdot M_2^*(X_1,\ldots, X_{m_1}), 
% \]
% $M=A,B,C$, and
% \begin{align*}
% K_{1,i} &= (1-X_0)\cdot  eq(i,X_1,\ldots, X_{m_1}), 
% \\
% K_{2,i} &= X_0\cdot eq^*(i,X_1,\ldots, X_{m_1}),
% \end{align*}
% for $i=0,\ldots, t$.
%The three claims \eqref{e:phase3:claim:vM} and \eqref{e:phase3:claim:inputs}  are combined into a single linear claim on $w$, 
Note that computing the matrix values \eqref{e:split:matrix:values} requires $O(n)$ field operations, where $n$ is the maximum number of non-zero entries in the R1CS matrices. 
While these values can be proven via the SPARK polynomial commitment scheme for sparse multilinears \cite{Spartan}, our main use case pivots for letting the verifier take the burden of matrix evaluation.

\begin{rem}
If WHIR is configured to consume only $m_1-k_r$ variables, for some small $k_r$, then the matrix polynomials in the final WHIR claim are specialized to $M(\vec\alpha, \vec\beta\|\:.\:)$, $M=A,B,C$, which are multilinears in the remaining $k_r$ variables. 
\end{rem}

% \paragraph{The combined inner product claim.} 

% The above mentioned linear combination of the claims \eqref{e:phase3:claim:vM} and \eqref{e:phase3:claim:inputs} is follows. 
% %Let us roll out the concrete form of the linearly combined constraints, in terms of the component polynomials:
% \begin{align}
% \label{e:combined:claim}
% v 
% &=\sum_{\vec x\in H_{m_1}} K_1(\vec \alpha, \vec x)\cdot w_1(\vec x)
% + \sum_{\vec x\in H_{m_1}} K_2^*(\vec\alpha,\vec x)\cdot w_2^*(\vec x), 
% \end{align}
% where
% \begin{align}
% K_1(\vec\alpha,\:.\:) &= A_1(\vec\alpha, \:.\:)  + \lambda B_1(\vec\alpha, \:.\:)  + \lambda^2 C_1(\vec\alpha, \:.\:)  
% + \lambda^3 \sum_{i=0}^{t-1} \lambda^i\cdot eq(i,\:.\:)
% \\
% K_2(\vec\alpha,\:.\:) &= A_2(\vec\alpha, \:.\:)  + \lambda B_2(\vec\alpha, \:.\:)  + \lambda^2 C_1(\vec\alpha, \:.\:)  
% + \lambda^{3+t} \sum_{i=0}^{t-1} \lambda^i\cdot eq(i,\:.\:)
% \end{align}
% and
% \begin{align*}
% v = v_A + \lambda\cdot v_B + \lambda^2\cdot v_C + \lambda^{3}\sum_{i=0}^{t} \lambda^{i}\cdot in_i + \lambda^{3+t}\sum_{i=0}^t \lambda^{i} \cdot\xi_i.
% \end{align*}
% Note that \eqref{e:combined:claim} is an inner product $v = \langle K, w\rangle_{H_{m_1 + 1}}$ of the combined witness polynomials $w = (1-X_0) w_1 + X_0 w_2^*$ with the combined kernel $K = (1-X_0) K_1(\alpha,\:.\:) + X_0 K_2^*(\alpha,\:.\:)$.





\subsection{Soundness.}
Without specifying the corresponding soundness notions,
let us informally state the soundness guarantees of the protocol, for proximity parameters up to the Johnson radius.
This means that $\theta$ as in the main relation \eqref{e:IOPP:main:R:0} satisfies
\[
\theta \in \left(0, 1 -\sqrt{1 - \delta_1}\right),
\]
where $\delta_1$ is the relative minimum distance of the base code $C_1 = C$. 
% At these distance $\theta$, the number of proximates to a 
The corresponding list size bound is
\[
\ell_1(\theta) =  \frac{1}{2 \sqrt{1-\delta_1}} \cdot \frac{1}{\eta},
\]
the combinatorial Johnson bound, with  $\eta =1 - \sqrt{1 - \delta_1} - \theta$.
WHIR with $r$ rounds involves the further Reed--Solomon codes $C_2,\ldots, C_r$, with each round having its own proximity parameter $\theta_i$, again assumed to be within the Johnson radius. 
(See Section \ref{s:WHIR}.)



% \[
% J(C_i) = 1 - \sqrt{1 - \delta(C_i)},
% \]
% of the base codes $C_1,\ldots, C_r$ used in WHIR. 
% By the Johnson bound, the number of proximate codewords at given distance is then
% \[
% \ell(\theta) < \frac{1}{2 \sqrt{1-\delta}} \cdot \frac{1}{\eta},
% \]
% where is the gap to the Johnson radius. 
% We do that 

\begin{thm}
\label{thm:main:soundness:informal}
For proximity parameters $\theta_i\in \left(0, 1 - \sqrt{1-\delta_i} \right)$, $i=1,2,\ldots, r$, where $\delta_i$ are the minimum distances of the Reed--Solomon codes $C_1$, \ldots, $C_r$.
Protocol \ref{prot:main} is a round-by-round knowledge sound interactive oracle proof for that $f_1\in \mathcal R_{\theta_1}$ as specified in \eqref{e:IOPP:main:R:0}, with the following overall soundness error for the three phases:
\begin{align}
\label{e:phase1:soundness:error}
\varepsilon_\text{phase 1} 
&\leq 
    % % \max\left( 
    % \binom{\ell_ 1}{2} \cdot \frac{2^{m_1}}{|F|} %\quad
    % + \ell_2\cdot \varepsilon_{stage 2} %\quad
    % + \binom{\ell_2}{2} \cdot \frac{2^{m_2}}{|F|},
    % % \max\right)
    \ell_1(\theta_1) \cdot \varepsilon_\text{R1CS} 
\\\intertext{%
where $\varepsilon_{R1CS}$ is the soundness error of the randomized AIR, 
}\label{e:phase2:soundness:error}
\varepsilon_\text{phase 2}  
&\leq \ell_1(\theta_1) \cdot \max\left\{\frac{m}{|F|}, \frac{3}{|F|}\right\},
\end{align}
and $\varepsilon_\text{phase 3} \leq 
%\ell_1\cdot \frac{3 + 2t - 1}{|F|} + 
\varepsilon_\text{WHIR}(\theta_1,\ldots, \theta_r)$ for the overall soundness error of WHIR with $r$ rounds on the inner product claims \eqref{e:phase3:claim:vM} and \eqref{e:phase3:claim:inputs}.
\end{thm}



% \begin{rem}
% Whenever the commitments $f_1$, $f_2$ belong to the same base code, and hence are absorbed in the same folding step of WHIR, the quadratic behavior in the list size is reduced to $\varepsilon_\text{phase 1} \leq \ell_1 \cdot \varepsilon_{stage 2}$ and $\varepsilon_\text{phase 2} \leq \ell_1\cdot \left(\frac{m}{|F|} + m\cdot \frac{3}{|F|}\right)$, since in this case, WHIR proves correlated agreement.
% See the discussion below.
% \end{rem}

% The soundness analysis behind Theorem \ref{thm:soundness:main:informal} is largely standard.
% Nevertheless we quickly discuss the soundness errors of the phases, and how lifting plays into it.
% (The role of lifting)
% \begin{equation}
% \label{e:IOPP:main:R:0}
% \mathcal R_\theta = \left\{
%     f_1\in \left(\mathbb F_q^{D}\right)^{e} \::\:
%     \begin{matrix}
%     \exists p(X) = \sum_{i=0}^{2^{m_1}- 1} c_i X^i,  %\in \mathbb F_q[X]^{<2^{m_1}}, 
%     d(p, f) \leq \theta,
%     \\
%     \text{ with }
%     \\
%     (c_i)  \in \mathscr R_{main} \cap \mathscr R_{aux},
%     \\
%     (c_1,\ldots,c_{t_1}) = \vec{in}
%     \end{matrix}
% \right\},
% \end{equation}

Let us roughly sketch the analysis behind Theorem \ref{thm:main:soundness:informal}.
First of all, the entire protocol aims to prove a proximate solution of the \textit{lifted} R1CS, specified by $A_1, B_1, C_1$ and the prefix-embedded second stage matrices $A_2^*, B_2^*, C_2^*$.
%which have the same number of columns as the first stage matrices.
%padded to the same number of columns as the first stage matrices.
% Unlike in \cite{lifted:FRI},
Unlike in certain univariate arithmetizations, see \cite{liftedFRI}, lifting by prefix-embedding does not change the circuit logic -- it only introduces variables that are never used.
The lifted system has the same soundness error as the original one, i.e. $\varepsilon_{R1CS^*} = \varepsilon_{R1CS}$, and from the point of view of $w_1$, it is equivalent to ask for a second stage response satisfying the lifted system, or a smaller witness satisfying the original system.\footnotemark
\footnotetext{%
In fact, a solution of the non-lifted system is simply obtained by truncation.
}
This is the reason why Phase 1 is a randomized reduction from the main claim $f_1\in \mathcal R_{\theta_1}$ to that the extended prover transcript $(f_1, f_2^*)$ belongs to 
\begin{align*}
\mathcal R_\text{phase 1} 
&=
\left\{
    \begin{matrix}
        f_1 \in \left(\mathbb F_q^{D}\right)^{e},% v_1\in F 
        \\
        f_2^* \in \left(\mathbb F_q^{D}\right)^{e}%, v_2\in F 
    \end{matrix}
    \::\:
    % \hspace*{-0.25cm}
    \begin{matrix}
        \exists (p_1(X),p_2(X)) \in \mathscr P_{m_1}^2 %\times \mathscr P_{m_1}
        \\
        d((p_1,p_2), (f_1,f_2^*)) \leq \theta_1,
        % \\
        % p_1(z_1) = v_1, p_2(z_2) = v_2.
        \\
        \text{and their coefficient vectors}
        \\
        (\vec c_1, \vec c_2) \text{ satisfy the lifted R1CS }
        \\
        \text{on inputs }\vec{in} , \vec\xi 
    \end{matrix}
\right\}.
\end{align*}
(Here, the distance refers to the interleaved representation of $p_1$ and $p_2$.)
%The latter will be ensured by Phase 2 and Phase 3 of the protocol.
Note, that the proximate $p_2(X)$ is not demanded of lifted form, in other words, it is not asked to stem from a $\theta_1$-proximate of the non-lifted word $f_2$ on $D'$.
(In fact, we do not rule out that $f_2$ is not as close to $C'$.)
A usual pigeon-hole argument on the number of $\theta_1$-proximates of $(f_1,f_2^*)$ bounds the soundness error of this reduction by
\[
\ell_1(\theta_1) \cdot \varepsilon_{R1CS^*} = \ell_1(\theta_1)\cdot \varepsilon_{R1CS},
\]
as stated in the theorem.
% Having a solution to the lifted R1CS is 

% It simply demands any proximate $(p_1(X), p_2(X))$ of $(f_1,f_2^*)$ which solves the \textit{lifted} R1CS,  however, truncation  automatically yields a solution of the non-lifted constraint system. 
% Thus we may use $\varepsilon_{R1CS}$ of the non-lifted system to argue the soundness error of Phase 1, scaled by $\ell_1(\theta_1)$ due to the usual pigeon-hole principle.
% Therefore, if the prover is able to respond on more than $\ell_1(\theta_1)\cdot \varepsilon_\text{R1CS}\cdot |F^t|$ many challenges $\vec\xi$ with a function $f_2^*$ so that $(f_1, f_2^*)\in \mathscr R_\text{phase 1}$, then by the pigeon-hole principle, at least one of the $\theta_1$-proximates of $f_1$ must have more than $\varepsilon_\text{R1CS}\cdot |F^t|$ many continuations which solves the non-lifted R1CS.

% The error for Phase 1 is the soundness error of randomized R1CS, scaled by the number of proximates of the committed functions. 
% Since both $f_1$ and $f_2$ are committed over the same base code, we view them as elements of the extended interleaved code
% \[
% (f_1,f_2)\in C_1^{2e},
% \]
% and $\ell_1$ is the list size bound at distance $\theta_1$, inherited from the base code $C_1$.
% WHIR will enforce correlated agreement between the components of the two, formally words of the interleaved code $C_1^{2e}$. 
% In general, $f_1$ and $f_2$ belong to different base codes, and we assume all combinations of the proximates,
% \[
% \ell_1 \cdot \ell_2
% \]
% as a crude upper bound for the proximates of the pair of commitments.
% Although we believe that a cross-domain analysis in the style of \cite{Stwo:whitepaper} can be extended to WHIR, thereby reducing this term, we do not pursue this direction here due to the substantial technical challenges involved.
% Our extension field $F$ provides enough security margin to handle quadratic (or even cubic) products of list sizes.
The same list size bound also scales the soundness errors of Phase 2, which is a reduction from $\mathcal R_\text{phase 1}$ to that the prover transcript $f_1,f_2^*$, extended by its proposed sumcheck polynomials $q_1(X),\ldots, q_m(X)$ and the final claims $v_A, v_B, v_C$, belongs to 
\begin{equation*}
\mathcal R_\text{phase 2} =
\left\{
    \begin{matrix}
        f_1 \in \left(\mathbb F_q^{D}\right)^{e}, %v_1\in F 
        \\
        f_2^* \in \left(\mathbb F_q^{D}\right)^{e_2}, %v_2\in F 
        \\
        q_1(X), \ldots, q_m(X) \in F[X]^{\leq 3}
        \\
        v_A, v_B, v_C \in F
    \end{matrix}
    \::\:
    \hspace*{-10mm}
    \begin{matrix}
        \exists (p_1(X),p_2(X)) \in \mathscr P_{m_1}^2
        \\
        d((p_1,p_2), (f_1,f_2^*)) \leq \theta_1,
        \\\text{and}
        \\
        0 = q_1(0) + q_1(1),
        \\
        q_i(\alpha_i) = q_{i+1}(0) + q_{i+1}(1),
        \\
        \text{for }i = 1,\ldots, m-1,
        \\
        q_m(\alpha_m) = (v_A \cdot v_B - v_C)\cdot eq(\vec\alpha, \vec\rho)
    \end{matrix}
\right\}.
\end{equation*}
The first term of the soundness error \eqref{e:phase2:soundness:error},
\[
\ell_1(\theta_1) \cdot \frac{m}{|F|}
\]
attributes the reduction from the domain identity to the inner product with $eq(\:.\:,\vec\rho)$.
%and corresponds to the probability that one of the proximates yields a non-zero constraint multilinear which evaluates to zero at the given point $\vec\rho$.
The other term, 
\[
\ell_1(\theta_1)\cdot \frac{3}{|F|},
\]
is the soundness errors of each of the $m$ rounds of the sumcheck protocol, on the cubic zero-check expression.

Finally, and without going into details, Phase 3 is a reduction from $\mathcal R_\text{phase 2}$ to the final verifier relation $\mathcal R_{V}$, which trades the correlated agreement claim
\[
\exists (p_1(X),p_2(X)) \in \mathscr P_{m_1}^2, \quad
        d((p_1,p_2), (f_1,f_2^*)) \leq \theta_1,
\]
for consistency checks on the WHIR sumcheck polynomials and the aggregated WHIR claim on the final polynomial revealed in plain. 
(The latter claim involves the evaluation claims on the matrix polynomials $M(\vec\alpha, \vec \beta)$.)
% \[
% \mathscr R_{V} =
% \left\{
%     \begin{matrix}
%         f_1 \in \left(\mathbb F_q^{D}\right)^{e}, %v_1\in F 
%         \\
%         f_2^* \in \left(\mathbb F_q^{D}\right)^{e}, %v_2\in F 
%         \\
%         q_1(X), \ldots, q_m(X) \in F[X]^{\leq 3},
%         \\
%         v_A, v_B, v_C \in F,
%         \\
%         q_{1,1}(X),\ldots, q_{1,\epsilon_1}(X) \in F[X]^{\leq 3},
%         \\
%         g_2 \in \left(\mathbb F_q^{D_{2}}\right)^{e_{2}}, u_2\in F, 
%         \\
%         \vdots
%         % \\
%         % q_{r-1,1}(X),\ldots, q_{r-1,\epsilon_1}(X) \in F[X]^{\leq 3},
%         % \\
%         % g_{r-1} \in \left(\mathbb F_q^{D_{r}}\right)^{e_{r}}, u_{r-1}\in F,
%         % \\
%         \\
%         q_{r,1}(X),\ldots, q_{r,\epsilon_r}(X) \in F[X]^{\leq 3},
%         \\
%         g_r \in \mathscr P_{k_r} 
%     \end{matrix}
%     \::\:
%     % \hspace*{-5mm}
%     \begin{matrix}
%         % \exists (p_1(X),p_2(X)) \in \mathscr L_{m_1}\big|_{z_1,v_1}\times \mathscr L_{m_2}\big|_{z_1,v_1}
%         % \\
%         % d(p_1, f_1) \leq \theta_1, d(p_2, f_2)\leq \theta_2
%         % % \\
%         % % p_1(z_1) = v_1, p_2(z_2) = v_2,
%         % \\\text{and}
%         \\
%         0 = q_1(0) + q_1(1),
%         \\
%         q_i(\alpha_i) = q_{i+1}(0) + q_{i+1}(1),
%         \\
%         \text{for }i = 1,\ldots, m-1,
%         \\
%         q_m(\alpha_m) = (v_A \cdot v_B - v_C)\cdot eq(\vec\alpha, \vec\rho)
%         \\\text{and}
%         \\
%             g_r \text{ satisfies the aggregated WHIR claim}
%         \\
%         \bar v = \langle K_r, g_r\rangle_{H_{k_r}}, 
%         \\
%         \text{which involves the specializations}
%         \\
%         M_i(\vec\alpha, \vec\beta \| \: .\:), 
%         \\ 
%         M=A,B,C, \: i=1,2
%     \end{matrix}
% \right\}.
% \]
Its soundness error is $\varepsilon_\text{WHIR}(\theta_1,\ldots, \theta_r)$, the overall error of WHIR with $r$ rounds on the five claims \eqref{e:phase3:claim:vM} and \eqref{e:phase3:claim:inputs}.
We describe the protocol in the next section.




\section{WHIR}
\label{s:WHIR}


% Given the commitment $g_0$ of a multil

% Given a function $g_0\in F^{D}$, and proximity parameter $\theta$, WHIR \cite{WHIR} proves proximity of $g_0$ to a polynomial $u_0(X)$ the multilinear representation of which,  $P_0(X_1,\ldots, X_{k_0})$, satisfies a linear claim of the form
% \[
% c_0 = \langle P_0, K_0\rangle_{H_{k_0}},
% \]

% Given an inner product claim $c_0 = \langle P_0, K_0\rangle_{H_{k_0}}$ between a multilinear $P_0(X_1,\ldots, X_m)$ with a public $K_0(X_1,\ldots, X_n)$, WHIR recursively transforms it into smaller claims 
% \[
% c_i = \langle P_i(\:.\:), K_i(\:.\:)\rangle_{H_{m-k_i}},
% \]
% where $P_i$ and $K_i$ are random linear combinations of the previous $P_{i-1}$ and $K_{i-1}$, respectively. 

% is a random specialization of the previous polynomial $P_{i-1}$, and $K_i$ an aggregate of the previous $K_{i-1}$ together with the their consistency spot checks.

For completeness we give a concise description of the WHIR \cite{WHIR} proof of proximity to constrained Reed--Solomon codes. 
Given a word defined on some set, $g_0:D\longrightarrow \mathbb F_q$, WHIR \cite{WHIR} is an interactive oracle proof of proximity to a univariate polynomial $u_P(X)$, the multivariate representation of which, say $P(X_1, \ldots, X_{k_0})$, satisfies a linear claim of the form
\begin{equation}
\label{e:WHIR:main:claim}
c_0 = \langle P_0, K_0\rangle_{H_{k_0}},
\end{equation}
where $K_0(X_1,\ldots, X_{k_0})$ is a publicly known (and often succinctly evaluable) multilinear. 
Similar to Basefold \cite{Basefold} and preceding works, e.g. \cite{Gemini,Hyperplonk} \cite{DualGemini:hackmd,Zeromorph}, the construction of WHIR rests on the connection between multilinear sumcheck specialization and the FRI-like folding of their univariate representations.

\subsection{Setup}
The mechanics of WHIR is in alignment with a specified sequence of univariate spaces
\begin{align}
\label{e:chain:function:spaces}
\mathscr P_{k_0} \supset \mathscr P_{k_1} \supset \ldots \supset \mathscr P_{k_r},
\end{align}
with $k_{0} > k_1 > \ldots > k_{r}\geq 0$. 
(These will correspond to the number of remaining variables after each round.)
As before, we commit polynomials in interleaved representation, using the linear isomorphism
\[
\mathscr P_{k_i} \simeq \mathscr P_{k_{i+1}}^{e_{i+1}}, \quad e_{i+1} = \dim{\mathscr P_{k_i}/\mathscr P{k_{i+1}}},
\]
by FFT decomposition \eqref{e:FFT:decomp}, for each $i=0,\ldots, r-1$.
The decomposed polynomials, elements from
\begin{equation}
\mathscr P_{k_1}^{e_1}, \mathscr P_{k_2}^{e_2}, \ldots, \mathscr P_{k_{r}}^{e_{r}},
\end{equation}
are then %evaluated over their own domain $D_i$, 
committed as words from the interleaved codes
\begin{equation}
 C_1^{e_1}, C_2^{e_2}, \ldots, C_r^{e_r},
\end{equation}
with specified base codes $C_{i} = \RS\left[F, D_{i}, 2^{k_{i}}\right]$ of (typically increasing) minimum distance.
Their evaluation domains $D_i$ are chosen as multiplicative cosets\footnotemark in $\mathbb F_q$, and are of size $|D_{i}|= 2^{n_{i}}$ for some $n_i > k_{i}$. 
\footnotetext{%
More generally, disjoint unions of multiplicative cosets of size $2^{k_i}$.
}
Each round further specifies $s_i\geq 1$ for the number of samples in the consistency check between adjacent oracles, as described below.


\subsection{The IOPP.}

In the following $\mathbb F_q$ is an FFT-friendly prime field, and $F$ a finite extension of it. 
The protocol starts with initial claim \eqref{e:WHIR:main:claim} and consists of $r$ rounds. 
Each round $i=1,\ldots,r$ reduces a claim $c_{i-1} = \langle P_{i-1}, K_{i-1}\rangle_{H_{k_{i-1}}}$ on the previous polynomial $P_{i-1}$ to a claim $c_{i} = \langle P_{i}, K_{i}\rangle_{H_{k_i}}$ on a new polynomial.

%typically cryptographically large.

% As described at the beginning of Section \ref{s:protocol:no:zk}, WHIR specifies a chain of univariate polynomial spaces
% \begin{align*}
% % \label{e:chain:function:spaces}
% \mathscr P_{k_0} \supset \mathscr P_{k_1} \supset \ldots \supset \mathscr P_{k_{r-1}} \supset \mathscr P_{k_r},
% \end{align*}
% with $k_{0} > k_1 > \ldots > k_{r}\geq 0$.
% A polynomial $p(X)$ from $\mathscr P_{k_{i-1}}$, $i = 1,\ldots, r$, is committed in decomposed manner as element $(p_1(X),\ldots, p_{e_i}(X))$ from  $\mathscr P_{k_i}^{e_i}$, where $e_i = \dim \mathscr P_{k_{i-1}} / \mathscr P_{k_i}$, according to the common FFT decomposition \eqref{e:FFT:decomp}.
% Thus, the codes met in the course of WHIR are the interleaved Reed--Solomon codes $C_{i}^{e_i}$, $i=1,\ldots, r$, with base code
% \[
% C_{i} = \RS[F,D_i,2^{k_i}], \quad i=1,\ldots, r,
% \]
% generated by the polynomials from $\mathscr P_{k_i}$, respectively.
% The evaluation domains $D_i$ are multiplicative subgroups of $\mathbb F_q$, of size $|D_i|=2^{n_i}$, for some $n_i > k_i$.


\begin{protocol}[WHIR]
\label{prot:WHIR}
The prover has a multilinear $P_{0}\in \mathbb F_q[X_1,\ldots, X_{k_0}]$, satisfying the inner-product claim
\[
c_0 = \langle P_{0}, K_0 \rangle_{H_{k_0}} 
\]
with some public multilinear $K_{0}\in F[X_1,\ldots, X_{k_0}]$.
The verifier has oracle access to $g_{0}\in (\mathbb F_q^D)^{e_1}$ the commitment of $P_0$'s univariate representation $u_{P_0}(X)\in \mathscr P_{k_0}$ as interleaved word from $C_1^{e_1}$.

Both prover and verifier engage in the following $r$ rounds, for $i=1,\ldots, r$, each consisting of three phases. 
\begin{enumerate}
    \item 
    Sumcheck for $r_i = k_{i-1} - k_i$ rounds, on the previous inner product claim between $P_{i-1}$ and $K_{i-1}$.
    In the course of these rounds, the prover has sent the quadratic refinements polynomials $q_{i,1}(X)$, \ldots $q_{i,r_i}(X)$, and the verifier has answered with $\vec\beta_i = (\beta_{i,1}, \ldots \beta_{1, r_i})\in F^{r_i}$.
    The prover commits to 
    \[
    P_{i} = P_{i-1}(\vec \beta_{i},\:.\:), 
    \]
    % the latter only when $k_i +1 \leq m_2$, 
    as the polynomial $u_{P_i}(X)\in \mathscr P_{k_i}$ represented by the interleaved word
    \[
    g_{i}  \in C_{i+1}^{e_{i+1}},
    \]
    and provides the verifier oracle access.
    (In the case $i=r$, the prover reveals the polynomial $u_{P_r}(X)$ in plain.)
    % and their next refinement polynomial
    % \begin{align*}
    % q_{k_i + 1}(X) = &\langle K_{i,1}(\:.\:,X,\beta_{k_i + 1}, \ldots,\beta_{m_1}), p_1(\:.\:,X,\beta_{k_i + 1}, \ldots, \beta_{m_1})\rangle
    % \\
    % &+ \langle K_2(\:.\:,X,\beta_{k_i + 1}, \ldots,\beta_{m_2}), 
    %  p_2(\:.\:,X,\beta_{k_i + 1}, \ldots,\beta_{m_2}) \rangle,
    % \end{align*}
    \item 
    Query phase. 
    \begin{enumerate}
    \item
    The verifier draws a DEEP query $x_{0} \sample F$, on which the prover answers with the value $v_{0}$ of the univariate representations $u_{P_i}(X)$ at $x_{0}$.  
    \item 
    The verifier queries $g_{i-1}$ from the previous round at $s_i>0$ random positions 
    \[
    x_{i,1}, \ldots, x_{i,s_i} \sample D_{i},
    \]
    and computes the expected values $v_{1}, \ldots, v_{s_i}$ for $u_{P_{i}}(X)$ at these points from the oracle answers, and $\vec\beta_i$.
    \end{enumerate}
    \item 
    Aggregation.
    The verifier samples $\gamma_i\sample F$ for combining the specialized sumcheck claim
    \begin{align*}
    q_{i, r_i}(\beta_{i,r_i}) &= \langle P_{i}, K_{i-1}(\vec \beta_{i},\:.\:) \rangle_{H_{k_i}} 
    \end{align*}
    and the evaluation claims
    \begin{align*}
    v_{j} &= \langle P_{i}, L_{x_i}\rangle_{H_{k_i}}, 
    \quad j = 0,\ldots, s_i,
    \end{align*}
    into a the new inner product claim
    $
    c_i = \langle P_{i}, K_{i}\rangle_{H_{k_i}} 
    $
    with the new $c_i$ and kernel $K_{i}$ as clarified below.
\end{enumerate}
If the consistency checks in all the sumcheck rounds hold, and if the final inner product claim $c_r = \langle P_r, K_r\rangle$ holds for the multilinear representation of the revealed polynomial $u_{P_r}(X)$, then the verifier accepts.
\end{protocol}

In the aggregation step, the combined kernel is
\begin{align*}
    K_{i} &=  K_{i-1,i,1}(\vec\beta_{i},X_1,\ldots, X_{k_i}) + \sum_{j=0}^{s_i} \gamma_i^{j+1} \cdot L_{x_{j}}(X_1,\ldots, X_{k_i}),
    \end{align*}
and $c_i = q_{i,r_i}(\beta_{i,r_i}) + \sum_{j=0}^{s_i} \gamma_i^{j+1}\cdot v_{j}$.
    

% \subsection{Combining several claims.}



\subsection{Soundness.}

% Let us quickle discuss the soundness properties of WHIR, for proximity parameters in the Johnson regime of the Reed--Solomon codes, 
% \[
% \theta_i \in \left(0, 1 -\sqrt{1-\delta_i}\right),
% \]
% where $\delta_i$ is the minimum distance of $C_i$, $i=1,\ldots, r$.

Every round of WHIR is a randomized reduction $\mathcal R_{i-1}\longrightarrow \mathcal R_i$ in the verifier randomnesses $\beta_{i,1},\ldots,\beta_{i,r_i}$, $x_{i,0},x_{i,1},\ldots, x_{i,s_i}$, and $\gamma_i$, where 
\begin{equation}
\label{e:WHIR:round:relation}
\mathcal R_i = \left\{
\begin{matrix}
g_i\in (\mathbb F_q^{D_{i}})^{e_{i}}
% \\
% q_0(X) \in F[X]^{\leq 2}
\end{matrix}
\::\:
\begin{matrix}
\exists p_i(X)\in \mathscr P_{k_i}, 
%p_i(X) = p{P_i}(X_1,\ldots, X_{k_i})%
\\
d(g_i,p_i)\leq \theta_i,
\\
\langle P_i, K_i\rangle_{H_{k_i}} = c_i 
% \\
% q_0(0) + q_0(1) = c_0
\end{matrix}
\right\}.
\end{equation} 
(Again, $P_i$ is the multilinear representation of $p_i$, and the distance refers to the interleaved representation of $p_i$.)
We use somewhat lazy notation here, hiding the dependency of $\mathcal R_i$ on the previous messages:
The inner product kernel $K_i$ incorporates the specialized sumcheck claim on the previous oracle $g_{i-1}$, its values at the $s_i$ in-domain queries, and the claimed value for $p_i$ at the single out-of-domain point $x_{i,0}$.
We note that in the case $i=r$, since $p_r(X)\in \mathscr P_{k_r}$ is revealed in plain, $g_r = p_r$, the distance condition in \eqref{e:WHIR:round:relation} becomes vacuous. 

The soundness analysis of the sumcheck phase, in parallel with the multilinear combinations of the interleaved components, is the most technically involved part. 
Its proof relies on (degree-$2$) curve decodability of Reed--Solomon codes, yielding correlated agreement under the additional constraints imposed by the inner products. 
% for \textit{constrained} Reed--Solomon codes \cite[Theorem 4]{Basefold:LDR}, a consequence of degree-2 curve decodability up to Johnson radius of $C_i$, 
% %which is a consequence of the more recent notion of line decodability \cite{goyalguruswami}.
% which we shall state in a simplified form in Appendix \ref{s:CAT}.
%
Without specifying the intermediate relations for each step of a WHIR round, we state its soundness properties as an (interactive oracle) reduction.
For a detailed analysis we refer to \cite[Section 6.4]{Basefold:LDR}, which generalizes in a straight-forward manner to arbitrary inner product constraints.
%of the reduction $\mathcal R_{i-1}\longrightarrow \mathcal R_i$.

\begin{thm}[WHIR Soundness]
For $\theta_i\in (0,1 -\sqrt{1-\delta_i})$, where $\delta_i$ is the minimum distance of the Reed--Solomon code $C_i$.
Then round $i$ in Protocol \ref{prot:WHIR} is round-by-round (knowledge) sound oracle reduction $\mathcal R_{i-1}\longrightarrow R_{i}$, with the following soundness errors for its steps.
\begin{enumerate}
    \item 
    Sumcheck phase:
    For each of the $r_i$ steps,
    \begin{align*}
       \varepsilon_{S,i} &= \varepsilon_{\text{CA}}(C_{i-1},2,\theta_{i-1}), 
    \end{align*}
    where $\varepsilon_\text{CA}(C_{i-1},2,\theta_{i-1})$ is the soundness error of the constrained correlated agreement theorem (Theorem \ref{thm:CAT:constrained}) for $C_{i-1}$, under degree-$2$ constraints, using proximity parameter $\theta_{i-1}$.
    \item 
    Query phase:
    \begin{align*}
        \varepsilon_{Q,i} = \frac{\ell_i(\theta_i)^2}{2} \cdot \frac{2^{n-k_{i}}}{|F\setminus D_i|} + (1 -\theta_{i-1})^{s_i},
    \end{align*}
    where $\ell_i(\theta_i)$ is the list size bound for $C_i$ at distance $\theta_i$.
    \item 
    Aggregation step:
    \begin{align*}
        \varepsilon_{A,i} = \ell_i(\theta_i) \cdot \frac{s_{i}}{|F|},
    \end{align*}
    again with the same list size bound as above.
\end{enumerate}
\end{thm}

\begin{rem}
For the soundness analysis, one may alternatively use the high-level notion of mutual correlated agreement \cite{WHIR}, a consequence of curve-decodability \cite{cryptoeprint:2025/2110,cryptoeprint:2025/2054}. 
\end{rem}
% Let us roughly discuss the soundness errors. 
% \begin{rem}
% Alternatively, one may use the stronger notion of mutual correlated agreement.
% \end{rem}

% Each step of the sumcheck is a randomized reduction $\mathcal R_{i-1}\stackrel{\lambda_i}{\longrightarrow} \mathcal R_i$, where
% \[
% \mathcal R_i = \left\{
% \begin{matrix}
% f_0\in F^{D_0},
% \\
% q_0, \ldots, q_{i-1} \in F[X]^{\leq 2}
% \end{matrix}
% \::\:
% \begin{matrix}
% q_{j-1}(\lambda_j) = q_j(0) + q_j(1), j=1,\ldots,i-1,
% \\
% \text{and}
% \\
% \exists p_i(X)\in \mathscr P_{n-i}, %P_i =\Phi^{-1}(p_i),
% \\
% d_{\pi^i(D_0)}(\mathsf{fold}_{\lambda_1,\ldots, \lambda_i}(f_0),p_i)\leq \theta_0,
% % \\\text{s.t.} 
% \\
% \langle P_i, R_0(\lambda_1,\ldots, \lambda_i,\:.\:)\rangle = q_{i-1}(\lambda_i),
% \end{matrix}
% \right\},
% \]
% for $i=1,\ldots, k-1$.


\section{Adding zero-knowledge}


\subsection{Randomizing $w_M(\vec\alpha)$.}

Recall that in the end of Phase \ref{i:main:phase:2} of Protocol \ref{prot:main}, the prover reveals the values
$w_A(\vec\alpha)$, $w_B(\vec\alpha)$, $w_C(\vec\alpha)$, which are subject to 
\begin{equation}
\label{e:zerocheck:last:check}
q_{m}(\alpha_m) = (w_A(\vec\alpha)\cdot w_B(\vec\alpha) - w_C(\vec \alpha))\cdot eq(\vec\alpha,\vec\rho),
\end{equation}
where $\alpha\in F^m$ are the random challenges collected over the rounds of the sumcheck.
In this section we discuss two different strategies to prevent these values from leaking private witness data.
%Both randomize either $w_1$ or $w_2$ ``on-chip'', by 
Both extend the witness (either $w_1$ or $w_2$) by random values, and add trivial gates to assure their propagation through the R1CS matrices. 



\subsubsection{Randomization by extension field elements.}

This is the canonical solution, where the openings over $F$ are randomized by random witnesses over $F$.
%This breaks with our current design decision keeping $w_2$ also over $\mathbb F_q$, so the solution would be only interesting if we decide to revert to $w_2$ over the extension field $F$.   
Note that it is sufficient to randomize $w_A(\vec\alpha)$ and $w_B(\vec\alpha)$; the value of $w_C(\vec\alpha)$ is then uniquely determined by \eqref{e:zerocheck:last:check}, as long as $eq(\vec\alpha, \vec\rho)\neq 0$. 

We allocate two private extension field variables $R_A, R_B$ in the second stage R1CS, and set their values to $r_A, r_B \sample F$. 
Their gates are 
\[
R_A \cdot 1 = R_A, \quad 1 \cdot R_B = R_B, 
%\quad 1 \cdot R_C = R_C,
\]
located at arbitrary row indices $\vec x_A, \vec x_B \in H_m$.
Then 
\begin{align}
\label{e:wA:randomized:extfield}
w_A(\vec Y) &= w_A'(\vec Y)  + r_A \cdot eq(\vec x_A,\vec Y),
\\
\label{e:wB:randomized:extfield}
w_B(\vec Y) &= w_B'(\vec Y) + r_B \cdot eq(\vec x_B, \vec Y) 
\end{align}
with $w_A'$, $w_B'$, not depending on the two random witnesses $r_A$, $r_B$.
 
% These two randomnesses are sufficient to hide the other witnesses, for most of the $\vec\alpha\in F^m$.
% $\vec x_1\in H_m$ corresponds to the row index of the gate $R_1\cdot 1 = R_1$ in $A_2$.

\begin{lem}
\label{lem:ext:field:randomization}
For any $\vec\rho\in F^m$.
Except for a fraction of $\varepsilon \leq \frac{3 m}{|F|}$ points $\vec\alpha\in F^m$, 
the distribution of $(w_A(\vec\alpha), w_B(\vec\alpha), w_C(\vec\alpha))$ is uniform on the set
\[
\left\{
(v_A, v_B, v_C) \in F^3 
\::\:
q_m(\alpha_m) = (v_A\cdot v_B - v_C)\cdot eq(\vec\alpha,\vec\rho)
\right\}.
\]
\end{lem}
\begin{proof}
The exceptional set is the union of the roots of the multilinears $eq(\vec x_A, \:.\:)$, $eq(\vec x_B, \:.\:)$  and $eq(\:.\:,\vec\rho)$.
For all other $\vec\alpha\in F^m$, the vector $(w_A(\vec\alpha), w_B(\vec\alpha))$ is uniformly distributed over $F^2$, since the coefficients of $r_A$ and $r_B$ do not vanish in \eqref{e:wA:randomized:extfield} and \eqref{e:wB:randomized:extfield}, and they uniquely determine $w_C(\vec\alpha)$ via the claim \eqref{e:zerocheck:last:check}.
$\Box$
\end{proof}

% We note that this already diverges from the univariate case, where local randomization of the witness propagates to \textit{all} points outside the witness domain. 

The bound in Lemma \ref{lem:ext:field:randomization} is essentially tight, as multivariate Lagrangians have exactly a fraction of $m/|F|$ zeros in $F^m$.
In other words, local randomization of the witness propagates only to \textit{almost all} points outside the witness domain.\footnotemark
\footnotetext{%
This diverges from univariate IOPs, where local randomization propagates to \textit{all} points outside the witness domain. 
}
% 
However, the fraction of exceptional points is of order $1/|F|$, and excluding them does not conflict our soundness goals.


For completeness we also cite an extension to multiple evaluation points:
\begin{lem}
\label{lem:multipoint:eval}
Choose $\Omega =\{\vec x_1, \ldots, \vec x_k\}$ a set of $k\geq 1$ distinct points on the Boolean hypercube $H_n$. 
Then, except with probability of at most $\varepsilon = \frac{k\cdot n}{|F|}$ in $\vec \alpha_1, \ldots, \vec \alpha_k \in F^n$, the multi-point evaluation matrix
\begin{equation*}
E_{\vec\alpha_1,\ldots,\vec\alpha_k} =
\begin{pmatrix}
eq(\vec x_1,\vec \alpha_1) & eq(\vec x_1, \vec\alpha_2) & \ldots & eq(\vec x_1, \vec\alpha_k)
\\
eq(\vec x_2,\vec \alpha_1) & eq(\vec x_2, \vec\alpha_2) & \ldots & eq(\vec x_2, \vec\alpha_k)
\\
\vdots & \vdots & &\vdots
\\
eq(\vec x_k,\vec \alpha_1) & eq(\vec x_k, \vec\alpha_2) & \ldots & eq(\vec x_k, \vec\alpha_k)
\end{pmatrix}
\end{equation*}
is of full rank.
\end{lem}

The proof shows that the determinant is a non-zero multilinear in the coordinates of $\vec\alpha_1, \ldots, \vec\alpha_k$.
We postpone it to the appendix.
% \begin{proof}
% The determinant of the matrix equals
% \[
% \sum_{\sigma\in S_k} L_{\vec x_{\sigma(1)}}(\vec\alpha_1) \cdot L_{\vec x_{\sigma(2)}}(\vec\alpha_2) \cdot \ldots \cdot L_{\vec x_{\sigma(k)}}(\vec\alpha_k) ,
% \]
% where $\sigma$ ranges over all permutations of $\{1,2,\ldots, k\}$, and therefore is a multilinear polynomial in the coordinates of $\vec\alpha_1$, \ldots, $\vec\alpha_k$.
% By Schwartz-Zippel, it can be zero for at most $k\cdot n$ points, unless it is identical zero.
% To exclude the latter, observe that each product in the determinant can be written as 
% \[
% L_{\vec x_{\sigma(1)}, \vec x_{\sigma(2)},\ldots, \vec x_{\sigma(k)}}(\vec \alpha_1,\ldots,\vec\alpha_k),
% \]
% the Lagrangian at $(\vec x_{\sigma(1)}, \vec x_{\sigma(2)},\ldots, \vec x_{\sigma(k)})$ in the $k\cdot n$-dimensional hypercube.
% Since the points $\vec x_1,\ldots, \vec x_k$ are distinct, so are the points in 
% \[
% \{(\vec x_{\sigma(1)}, \vec x_{\sigma(2)},\ldots, \vec x_{\sigma(k)}) \::\: \sigma\in S_k \},
% \]
% and we conclude from linear independence of Lagrangians that their sum cannot be zero.
% $\Box$
% \end{proof}




\subsubsection{Randomization by basefield elements.}
%In the light of Lemma \ref{lem:multipoint:eval} 
% As we would like to keep the witness over the basefield $\mathbb F_q$

It is tempting to believe that the univariate strategy from \cite{cryptoeprint:2024/1037}, which randomizes an extension field opening by extension-field-degree many basefield values, carries over to the multilinear world in a straight-forward manner.
However, the best generic result we found is the following one, which has a way too large bound for the exceptional set.

\begin{lem}
\label{lem:Lagrange:random:eval}
Let $F$ be a degree $e$ extension of a prime field $\mathbb F_q$, and $\Omega =\{\vec x_1, \ldots, \vec x_e\}$ a set of $e$ distinct points on the Boolean hypercube $H_m\subset \mathbb F_q^m$. 
Then, except for a fraction of at most $\varepsilon = \frac{m}{q-1}$ points $\vec \alpha \in F^m$, the elements
\begin{equation*}
eq(\vec x_1, \vec \alpha), \ldots, eq(\vec x_e, \vec \alpha) \in F
%= eq(\vec \rho,\vec x_i)    
\end{equation*}
are $\mathbb F_q$-linear independent. 
\end{lem}


The proof of the lemma relies on the well-known Moore determinant as a distinguisher for $\mathbb F_q$-linear independence.
However, we require the property also over $K$ the field of rational functions in several variables, and hence state it in full generality.
%on , which is as follows.
\begin{lem}
\label{lem:goss}
Let $K$ be a (possibly infinite) field of characteristic $q<\infty$, and let $f_1,\ldots,f_k\in K$, $k\geq 1$.
Then $f_1,\ldots, f_k$ are $\mathbb F_q$-linear independent if and only if the determinant of the Moore matrix
\[
% \Delta(f_1,\ldots, f_k) = \left|
M(f_1,\ldots, f_k) =
\begin{pmatrix}
f_1 & f_2 & \ldots & f_k
\\
f_1^q & f_2^q & \ldots & f_k^q
\\
\vdots &\vdots & \ddots & \vdots
\\
f_q^{q^{k-1}} & f_2^{q^{k-1}} & \ldots & f_k^{q^{k-1}}
\end{pmatrix}
% \right| \neq 0
\]
is a non-zero element in $K$. 
\end{lem}

For a proof of Lemma \ref{lem:goss}, we refer the reader to \cite[Lemma 1.3.3]{Goss96}.


\begin{proof}[of Lemma \ref{lem:Lagrange:random:eval}]
Since the points $\vec x_1,\ldots, \vec x_n$ are distinct, the Lagrangians $eq(\vec x_i,\vec Y)$, $i=1,\ldots,e$, are $\mathbb F_q$-linear independent elements of $K= \mathbb F_p(\vec Y)$, the rational function field in $m$ variables. % the coordinates of $\vec Y$.  
By Lemma \ref{lem:goss}, their Moore determinant 
\[
\det M(eq(\vec x_1, \vec Y),\ldots, eq(\vec x_e, \vec Y))\neq 0,
\]
as an element of $K$.
Moreover, since all matrix entries are polynomials, we conclude that the Moore determinant
is a non-zero \textit{polynomial}. 
% (All entries of the Moore matrix are polynomials.)
%
% Let us determine its total degree. 
The $(i,j)$-th entry in the Moore matrix is
% \[
% L_{x_j}(Y_1,\ldots, Y_n)^{p^{i}} = L_{x_j}(Y_1^{p^i},\ldots, Y_n^{p^i}),
% \]
has total degree at most $m\cdot q^{i-1}$, and the total degree the determinant is at most
\begin{align*}
m + m q + \ldots + m  q^{e-1} = m \cdot \frac{q^e - 1}{q - 1}.
\end{align*}
By Schwartz-Zippel, the fraction of $\vec\alpha\in F^m$ for which the determinant evaluates zero
is at most
\[
\frac{m}{q-1} \cdot \frac{q^e -1}{q^e} < \frac{m}{q-1}.
\]
For all other $\vec\alpha$ the determinant is non-zero, and another application of Lemma \ref{lem:goss} yields that $eq(\vec x_1,\vec\alpha),\ldots, eq(\vec x_e, \vec \alpha)$ are $\mathbb F_p$-linear independent elements of $F$.
$\Box$
\end{proof}

It is hard to tell whether the Schwartz--Zippel bound in Lemma \ref{lem:Lagrange:random:eval} is tight. 
A high polynomial degree is not a guarantee for a large set of zeros. 
% For example, 
% \[
% eq(\vec x_i, \vec Y)^{q^{i-1}} = eq(x_{i,1},\ldots, x_{i,m}, Y_1^{q^{i-1}},\ldots, Y_m^{q^{i-1}})
% \]
% is of total degree $m\cdot q^{i-1}$. 
% Its zeros are the roots of $Y_j^{q^{i-1}} - x_{i,j}$, but in $F$ there is only one solution, $Y_j = x_{i,j}$, since is no $q^{i-1}$-th root of unity, except $1$.\footnotemark
% \footnotetext{%
% This follows from $(q^{i-1},q^{e} - 1) = 1$. 
% }
% Thus the polynomials has the same number of roots as $eq(\vec x_i, \vec Y)$. 
We tried shortly with GPT-5.6, but did not achieve significant improvements. 
However, it proposed the following work around, which explicitly uses the multiplicative structure of the Lagrangians. 

\paragraph{Blowing up soundness.}
Instead of a single set $\Omega$ of size $e$, we choose several sets
\[
\Omega_1, \ldots \Omega_s \subset H_m,
\]
in a very particular manner, so that their exceptional sets are statistically independent.
(The latter will be possible because of the tensor structure of the Lagrangians.)
Each $\Omega_j$ is a set of $e$ points contained in a sub-cube of small dimension 
\[
t = \ceil{\log_2 e},
\]
with its free part touching the variables $Y_{i_j(1)}, \ldots, Y_{i_j(t)}$ only, meaning that its points do not vary outside these $t$ coordinates.  
%We call the variables $Y_{i_j(1)}, \ldots, Y_{i_j(s)}$ \textit{free variables}, and the remaining $m-s$ variables the \textit{fixed variables} of $\Omega_j$.
We further require that the free variables are disjoint between the sets, i.e.
\[
\big\{Y_{i_j(1)},\ldots, Y_{i_j(t)} \big\} \cap \big\{Y_{i_{j'}(1)},\ldots, Y_{i_{j'}(t)} \big\} = \varnothing,
\]
whenever $j\neq j'$.
The crux is that for each such set $\Omega_j$, we can drop its fixed variables in the proof of Lemma \ref{lem:Lagrange:random:eval}, as long as their $eq$-terms are not zero, since a common non-zero factor does not change $\mathbb F_q$-linear independence. 
The exceptional set for $\Omega_j$ is only in the variables $Y_{i_j(1)}, \ldots, Y_{i_j(t)}$, independent of the other variables, and we arrive at the following statement.

\begin{lem}
With $\Omega_1,\ldots, \Omega_s \subset H_{m}$ as above.
Then except for an $\varepsilon$-fraction of points $\vec\alpha\in F^m$, where
\begin{equation}
\label{e:basefield:randomization:error}
\varepsilon < s\cdot \frac{m-t}{|F|} + \left(\frac{t}{q-1}\right)^s,
\end{equation}
at least one of the sets $\Omega = \Omega_j$ is such that its values $eq(\vec x, \vec\alpha)$, $\vec x\in \Omega$, are $\mathbb F_q$-linearly independent.
\end{lem}
\begin{proof}
The first term bounds the the probability that $\vec\alpha\in F^m$ hits one of the $m-t$ fixed variables of $\Omega_j$, times $e$ for the union bound over all sets.
For all other $\vec\alpha$, the exceptional sets in Lemma \ref{lem:Lagrange:random:eval} are statistically independent, and of relative size $t/(q-1)$ each.
$\Box$
\end{proof}

In our application we would use $s = e$, to make the second term in \eqref{e:basefield:randomization:error} of order $O(1/|F|)$. 
That is,  we would require $e^2$ random witnesses over $\mathbb F_q$ to randomize $w_A(\alpha)$, and another $e^2$ random witnesses for $w_B(\alpha)$. 
They can be placed at arbitrary positions in the witness, yet their constraints (as before of the form $R_j \cdot 1 = R_j$, or $1\cdot R_j = R_j$) have to respect the positioning described above. 

%The main disadvantage is that we cannot place them at arbitrary positions. 



\subsection{Zero-check masking}
\label{s:zk:zerocheck}

The zero-check, Phase \ref{i:main:phase:2} in Protocol \ref{prot:main}, is turned into zero-knowledge by using the sumcheck randomization from Libra \cite{cryptoeprint:2019/317}.
The prover commits to the univariate polynomials 
\begin{equation}
\label{e:zk:zerocheck:randomizers}
h_0\sample F, \quad h_1(X),\ldots, h_m(X)\sample F[X]^{<d},
\end{equation}
where $d=3$ is the degree of the sumcheck \eqref{e:R1CS:zerocheck}.
These polynomials define the sparse mask
\begin{equation}
\label{e:zk:zerocheck:sparse:mask}
h(X_1,\ldots, X_n) = h_0 + X_1\cdot h_1(X_1) + \ldots + X_m\cdot h_n(X_m),
\end{equation}
which is overlaid with the hypercube sum \eqref{e:R1CS:zerocheck}. 
The prover claims the sum $c' = \sum_{\vec x \in H_m} h_m(X_m)$, the computation of which is almost for free via the succinct representation
\begin{align*}
\frac{1}{2^m}\sum_{\vec x\in H_m} h(\vec x) 
% = 2^{m} h_0 + 2^{m-1} h_1(1) + 2\cdot \sum_{\vec x\in H_{m-1}} X_2 h_2(x_2) + \ldots + x_m\cdot h_m(x_m) 
% \\
%  = 2^{m} h_0 + 2^{m-1} h_1(1) + 2\cdot (2^{m-2} h_2(1) + 2 \cdot \sum_{\vec x H_{m-2}} ...)
%  \\
=  h_0 + \frac{1}{2}\left(
h_1(1) + \ldots + h_m(1)
\right).
\end{align*}
%It receives a random $\alpha\sample F$ from the verifier, and 
Then both run the sumcheck on the combined claim
\begin{equation}
\label{e:zk:zerocheck:combined:claim}
c_0 = \alpha \cdot c + c' = \sum_{\vec x\in H_n} \alpha\cdot G(\vec x) + h,
\end{equation}
where $\alpha\sample F$ is random, and  $c = \sum_{x\in H_n} G(x_1,\ldots, x_n)$ is the target claim. 
(In our case $G$ is be the multivariate zero-check term, and $c=0$.)

%
The sparse mask carries sufficient entropy to perfectly hide the witness in the sumcheck polynomials
\begin{equation*}
\hat q_k(X) = \sum_{\vec x\in H_{n-k}} \alpha\cdot G(\alpha_1,\ldots, \alpha_{k-1},X,\vec x) +  h(\alpha_1,\ldots, \alpha_{k-1},X,\vec x),
\end{equation*}
for $k=1,\ldots,n$, as the following lemma shows.

\begin{lem}
For any $\alpha\in F $, $(\alpha_1,\ldots, \alpha_n)\in F^n$. 
The distribution of the sumcheck prover messages $c_0, \hat q_1(X),\ldots, \hat q_n(X)$ is uniform on the set
\[
\mathcal R = \left\{
\begin{matrix}
c_0\in F, \hat q_1,\ldots,\hat q_n\in F[X]^{\leq d}
\end{matrix}
\::\:
\begin{matrix}
c_0 = \hat q_1(0) + \hat q_1(1),
\\
\hat q_{k}(\alpha_{k}) = \hat q_{k+1}(0) + \hat q_{k+1}(1),
\\
k=1,\ldots n-1
\end{matrix}
\right\},
\]
independent of $G$.
\end{lem}

For the sake of completeness, we give a proof the lemma.
\begin{proof}
Fix $G$, $\alpha\in F$ and $(\alpha_1,\ldots,\alpha_n)\in F^n$, and choose any set $S\subset F$ of $d$ distinct points, containing the element $1$, but not $0$.
Using the consistency equations it is easily seen that the linear mapping, which sends $(c_0,\hat q_1(X),\ldots, \hat q_n(X))$ to the values   
\begin{equation}
\label{e:zk:zerocheck:polys:by:value}
(\vec v_1,\ldots, \vec v_n) = (\hat q_1|_{S\cup\{0\}},\hat q_1|_S, \ldots, \hat q_n|_S) 
\end{equation}
is a bijection between $\mathcal R$ and $F^{dn + 1}$. 
We leave this to the reader.\footnotemark
\footnotetext{%
The values $\vec v_1$ uniquely determine the polynomial $\hat q_1(X)\in F[X]^{\leq d}$ by degree, and with it $c_0 = \hat q_1(0) + \hat q_1 (1)$.
Likewise, $\hat q_1(\alpha_1)$ together with the values $\vec v_2$ uniquely determine $\hat q_2(X)\in F[X]^{\leq d}$ satisfying $\hat q_1(\alpha_1) = \hat q_2(0) + \hat q_2(1)$, since the letter determines the value of $\hat q_2(X)$ at the additional point $0$.
And so on.
}
To prove the lemma, it suffices to show that the values \eqref{e:zk:zerocheck:polys:by:value} are in one-to-one correspondence with $h_0\in F, h_1,\ldots, h_n\in F[X]^{<d}$.
But this is evident from the concrete form
\begin{multline*}
\hat q_{k}(X) = \sum_{\vec x \in H_{n-k}} \alpha \cdot G(\alpha_1,\ldots, \alpha_{k-1}, X,\vec x) 
\\
+ 2^{n-k} (h_0 + \alpha_1 h_1(\alpha_1) +\ldots +\alpha_{k-1} h_{k-1}(\alpha_{k-1})+ X h_{k}(X)) 
\\+ \sum_{\vec x \in H_{n-k}} x_{k+1} h_{k+1}(x_{k+1}) + \ldots + x_n 
h_n(x_n),
\end{multline*}
in which the other $h_0,\ldots, h_{k-1}, h_{k+1},\ldots, h_n$ impact only the constant part of the polynomial.
(The representation partially collapses in the two edge cases.
For $k=1$ there are no terms except $h_0$ preceding $X h_1(X)$, and for $k=n$, there are no terms following $X h_n(X)$.)
$\Box$
\end{proof}


The sumcheck reduces the claim \eqref{e:zk:zerocheck:combined:claim} to that 
\begin{align*}
\hat q_m(\alpha_m) = G(\alpha_1,\ldots, \alpha_m) + h_1(\alpha_1) + \ldots + h_m(\alpha_m),
\end{align*}
where
\begin{align*}
G(\vec\alpha) =  (w_A(\vec\alpha)\cdot w_B(\vec\alpha) - w_C(\vec\alpha))\cdot eq(\vec\alpha,\vec\rho),    
\end{align*}
for which the prover sends the values $w_A(\vec\alpha)$, $w_B(\vec\alpha)$, $w_C(\vec\alpha)$. 
The complementary term 
\begin{equation}
\label{e:zk:zerocheck:mask:opening}
h_1(\alpha_1) + \ldots + h_m(\alpha_m) = \hat q_m(\alpha_m) - G(\vec\alpha)
\end{equation}
can be proven in zero-knowledge by a separate protocol.

\begin{rem}
In our current implementation, we append the $dm+1$ coefficients\footnotemark for $h_0,h_1,\ldots, h_m$ to the first stage witness $\vec w_1$, and use zk-WHIR with the additional constraint
\begin{equation}
\hat q_m(\alpha_m) - G(\vec\alpha) = \sum_{\vec x\in H_{m_1}} I_\alpha(\vec x)\cdot  w_1(\vec x),
\end{equation}
where $I_\alpha$ is the inner product kernel for $h_1(\alpha_1) + \ldots + h_m(\alpha_m)$.
This requires the witness to be sufficiently below a power two. 
See the complete protocol in Section \ref{s:complete:protocol}.
\footnotetext{%
Alternatively, one may directly use the representation by their values \eqref{e:zk:zerocheck:polys:by:value}.} 
\end{rem}

Let us quickly sketch the transcript simulator for the zero-check.
The prover chooses arbitrary random polynomials $w_1(X_1,\ldots, X_{m_1})$  and $w_2(X_1,\ldots, X_{m_2})$, with their image polynomials $w_A, w_B, w_C$ not satisfying the zero-check hypercube sum for random $\vec\rho\sample F$.
Then it samples $h_0,h_1,\ldots, h_n$ and $\lambda\neq 0$ and computes the hypercube sum
\begin{equation}
\label{e:sumcheck:randomized:G}
c_0 = \sum G + \lambda h.
\end{equation}
It then runs the sumcheck protocol on the valid claim \eqref{e:sumcheck:randomized:G}.
The (wrong) claim $c'$ is chosen adaptively so that $\alpha \cdot 0 + c' = c_0$. 
For details, see Section \ref{}.

% For on-domain randomization we rely on the following Lemma. 
% \ulrich{We need an adaption to multivariates of bounded individual degree, if we reveal cross-terms}
% \begin{lem}
% Fix any $k$ distinct points $\vec x_1,\ldots, \vec x_k\in H_n$. 
% Then 
% \[
% \Pr_{\vec y_1,\ldots, \vec y_k\sample F^n}\left[
% \det \begin{pmatrix}
% eq(\vec y_i,\vec x_j)
% \end{pmatrix}
% = 0
% \right] 
% \leq \frac{k\cdot n}{|F|}.
% \]
% \end{lem}
% \begin{proof}
% The determinant is multilinear in the coordinates of $\vec y_1, \ldots, \vec y_k$.
% $\Box$
% \end{proof}


\subsection{Masked commitments.}
\label{s:masked:commitments}

To avoid witness information leakage from commitment openings, we use a slightly different approach than \cite{cryptoeprint:2026/391}.
Instead of committing to a zero-knowledge encoding of witness,
we overlay a witness polynomial with a sparse random mask without  degree increase:
Given witness $p \in \mathscr P_{m}$ and random $\msk\sample F[X]^{< q_i}$, for some small $q_i < 2^m$, we commit to
\begin{equation}
\hat p = p + \msk, %\in \mathscr P_m,
\end{equation}
which is again from $\mathscr P_m$.
(This secures the witness for up to $q_i$ openings.)
Obviously this does not bind the witness, unless we also commit to $\msk$, which is required for ZO0K anyways (see Section \ref{s:zook}).

For practical reasons, we \textit{jointly} commit all\footnotemark masks required by the protocol in one shot.
\footnotetext{%
These are the masks for the two witness polynomials $w_1$ and $w_2$, and the further masks $\msk_{1},\ldots, \msk_{r}$, $h_1,\ldots, h_r$ required for ZO0K.
}%
We concatenate their coefficient vectors, append $q_{zk}$ additional random coefficients, and encode the resulting vector as a single word from 
\begin{equation}
\label{e:C:zk}
C_{zk} = \RS[F,D_{zk}, d_{zk}].
\end{equation}
Here the degree bound $d_{zk}$ is at least $q_{zk}$ larger than the total of all mask lengths, and can be chosen again a power of two, if this helps encoding performance.
(The choice of $q_{zk}$ depends on the number of openings in the final proof of proximity, Protocol \ref{prot:zook:final:opening}.)

The rationale behind masked commitments is the following.
Adding the small mask $\msk$ in coefficient representation is a local modification of the witness, and the 
subsequent encoding of $\hat p$ is as costly as that of $p$. 
In contrast, even a tiny degree increase of the witness polynomial would lead to an $O(m)$ overhead in the encoding step, where $m$ is the degree of $p$.




\subsection{WHIR in zero-knowledge}
\label{s:zook}


We describe an adaptation of ZO0K \cite{cryptoeprint:2026/391} to our technique for randomizing commitments as described in the previous section.
The changes to the original protocol are minimal, and based on the fact that we can now also view masked polynomials as multilinears in the sumcheck claim.
See Remark \ref{rem:zook:pi:instead:pi:hat} for a one-to-one variant of \cite{cryptoeprint:2026/391}.


% \subsubsection{ZO0K.}
%ZO0K instead uses throughout sparse masks for randomization, and is therefore particularly prover friendly. 
The sumcheck rounds are randomized using the sparse Libra mask explained in Section \ref{s:zk:zerocheck}, and information leakage via commitment openings is mitigated through the commitment masks as described above.
In this multi-polynomial regime, ZO0K considers WHIR as a sumcheck prover, which gradually reduces composite claims of the form 
\begin{align}
\label{e:zook:claim}
c_{i}  &= \left\langle \hat p_{i}, K^{(i)} \right\rangle_{H_{k_{i}}} 
+ \sum_{j=0}^{i} \Gamma_j^{(i)}(\msk_j) + \sum_{j=1}^{i} \Lambda_{j}^{(i)} (h_{j}),
\end{align}
involving a costly part, the inner product $\left\langle \hat p_{i}, K^{(i)} \right\rangle$ of a large (masked) witness polynomial $\hat p_i$ with a public kernel $K^{(i)}$, and a cheap part, certain public linear functionals $\Lambda_j^{(i)}$, $\Gamma_j^{(i)}$ on the (small) masks committed by the prover so far:\footnotemark
\footnotetext{%
The $\Lambda_j^{(i)}(h_j)$ and $\Gamma_j^{(i)}(\msk_i)$ can be again represented as inner products. 
However, because of their fundamentally different treatment in the protocol, we write them as linear functionals.
}
The sparse sumcheck masks $h_1,\ldots, h_i$ and the witness masks $\msk_0,\ldots,\msk_i$.
The cheap part is not further reduced by the sumcheck logic; it is treated as constant function 
% \[
% \frac{1}{2^{k_i}}\cdot \sum_{j=1}^{i-1} \Lambda_{j}^{(i)} (h_{j}) + \Gamma_j^{(i)}(\msk_j)
% \]
on the hypercube $H_{k_i}$, which is passed through the rounds in a weighted form.
%leading to an simple update law for the linear functionals. 
The costly part is reduced in the usual sense, by folding. 
However, not before a fresh sumcheck mask $h_{i+1}$ enters the sum.
In this manner, the number of masks in the claim \eqref{e:zook:claim} grow from round to round, but the costly part shrinks.

After the final round, the remaining claim, a linear functional in $\hat p_r$, $h_1,\ldots, h_r$, $\msk_0,\ldots, \msk_{r}$, is proven in zero-knowledge by an elementary proof of proximity, Protocol \ref{prot:zook:final:opening} below.




% Recall that WHIR in full generality is a sumcheck prover, and works for arbitrary algebraic expressions $G(f_1,\ldots, f_m)$ in several multilinear polynomials $f_1$, \ldots $f_m$. 
% Every now and then it accompanies the sumcheck with new proposals for the multilinears in the specialized claim, and the claim of proximity to a sumcheck consistent ensemble of polynomials is reduced to a fresh claim on the new ensemble of smaller codewords.
% However there is no strict need to provide new proposals for every multilinear in the claim. 
% If a multilinear is small, e.g. depending only one or two variables or has a succinct representation, then the verifier can compute its specialization on its own, without sacrificing efficiency.

% This is the view point we take on ZO0K:
% In each round the prover replaces the (one and only) large polynomial in the claim by a smaller one. 
% On the other hand it adds few more small commitments required for zero-knowledge.
% The number of committed words grows, but the total size of the multilinears becomes smaller and smaller.


% (Here, we again use the same notation for the composed polynomial and its interleaved representation.)
% This is a linear constraint on the committed codewords $\hat p_0$ and $\hat\msk_0$.
 %(Their final specialized claim will be proven on their initial unfolded form.)
% Once the polynomial becomes 

\subsubsection{The required masks.}
As pointed out in Section \ref{s:masked:commitments}, we assume that the prover has committed beforehand to all the masks, as a single word from $C_\text{zk}$.
Concretely, each sumcheck mask $h_i$ is committed as coefficient vector 
\begin{equation}
\label{e:zook:sumcheck:mask:coeffs}
h_{i,0}\| h_{i,1} \| \ldots\| h_{i,r_i} \in F^{2 r_i + 1}
\end{equation}
from $h_{i,0}\sample F$, and $h_{i,j}\sample F[X]^{<2}$, appended by the coefficients of the commitment masks \begin{align}
\label{e:zook:commitment:mask:coeffs}
msk_i \in F^{q_i},  \quad i=0,\ldots,r,
\end{align}
from $\msk_i \sample F[X]^{< q_i}$, where $msk_0$ can be taken from the subfield $\mathbb F_q$, assuming that the first witness polynomial $p_0$ is over $\mathbb F_q$.
(In our application, our $p_0$ is the composite witness polynomial, both of which have their own basefield mask.)
The number of coefficients is chosen according to how often the commitment of $p_i + \msk_i$ is queried in the course of WHIR:
\begin{align}
q_0 &= s_{1}, 
\\
q_i &= s_{i+1} + 1 \text{ for } 1\leq i\leq r-1,
\\\intertext{assuming a single DEEP query in each of the rounds, and}
q_r &= s_\text{zk},
\end{align}
the number of queries on $\hat p_r$ in the final opening protocol, Protocol \ref{prot:zook:final:opening}.
The concatenated coefficient vectors \eqref{e:zook:sumcheck:mask:coeffs} and \eqref{e:zook:commitment:mask:coeffs} are encoded in zero-knowlege manner, by appending to the message another $q_{zk}$ random coefficients
\begin{equation}
\label{e:zook:mask:blinder:coeffs}
msk_{zk} \in F^{q^{zk}},
\end{equation}
before the Reed--Solomon encoding. 
That is, $C_{zk}$ encodes   
\[
h_1\|\ldots \|h_r \| msk_0 \|\ldots \| msk_r \| msk_{zk}
\]
parsed as coefficient vector from $F^{r + k_0 - k_r + q_0 + \ldots + q_r + q_{zk}}$ as Reed--Solomon codeword over $F$ with evaluation domain $D_{zk}$.

\subsubsection{Initialization.}
The initial claim is on the unmasked polynomial $p_0 = \hat p_0 - \msk_0 \in \mathscr P_{k_1}^{e_1}(\mathbb F_q)$, and of the form
\[
c_0 = \langle \hat p_0 - \msk_0, K \rangle_{H_{k_0}},
\]
where the inner product kernel $K$ is public.
This is of the form \eqref{e:zook:claim} with 
\[
K^{(0)} = K, \quad \Lambda_0^{(0)}(\msk_0) = \langle K, \msk_0\rangle_{H_{k_0}}.
\]
The latter can be easily computed as a linear functional on the coefficient vector $msk_0\in\mathbb F_q^{q_0}$:
The multilinear representation of $\msk_0\in F[X]^{< q_0}$ is non-zero only on the first $q_0$ points of the hypercube and thus
\[
\langle K, \msk_0\rangle_{H_{k_0}} = \sum_{j=0}^{q_0-1} K(j)\cdot msk_{0,j},
\]
in the index notation of the hypercube. 

\subsubsection{Modifications of a WHIR round.}

Let us describe the changes of round $i$ in Protocol \ref{prot:WHIR}, $i=1,\ldots, r$.
At the beginning of the round, the prover has already committed to 
$\hat p_{i-1}$, $\msk_0,\ldots, \msk_{i-1}$, $h_1,\ldots h_{i-1}$, satisfying the linear constraint
\begin{align}
\label{e:zook:prev:claim}
c_{i-1}  &= \left\langle \hat p_{i-1}, K^{(i-1)} \right\rangle_{H_{k_{i-1}}} 
 + \sum_{j=0}^{i-1} \Gamma_j^{(i-1)}(\msk_j)
+ \sum_{j=1}^{i-1} \Lambda_{j}^{(i-1)} (h_{j})
% + \sum_{j=1}^{i-1} \sum_{k=0}^{r_j} \Lambda_{i,j}^{(i-1)} (h_{j,k})
%     + \sum_{j=1}^{i-1} \Gamma_i^{(i-1)}(\msk_j)
\end{align}
We think of the claim as $c_{i-1} = \sum_{\vec x\in H_{k_{i-1}}} G_{i-1}(\vec x)$ for the quadratic expression
\begin{align*}
G_{i-1}(\vec x) = \hat p_{i-1}(\vec x) \cdot K^{(i-1)}(\vec x) - 2^{-k_{i-1}}\cdot B_{i-1},
\end{align*}
where $B_{i-1}$ is the sum of the mask functionals in \eqref{e:zook:prev:claim}.


\paragraph{Step 1: Masking phase.}
Identical to the sumcheck randomization in Section \ref{s:zk:zerocheck}.
The prover takes the masks $h_{i,0}\in F$, $h_{i,1},\ldots,h_{i,r_i} \in F[X]^{<2}$ it committed before running the protocol. 
They define the sumcheck mask
% \ulrich{mayb /}
\begin{equation}
\label{e:zook:sumcheck:mask}
h_i(X_1,\ldots, X_{r_i}) = h_{i,0} + X_1 \cdot h_{i,1}(X_1) + \ldots + X_{r_i} \cdot h_{i,r_{i}}(X_{r_i}),
\end{equation}
in as many variables as the sumcheck phase of the WHIR round consumes, $r_i = k_i - k_{i-1}$.
% in the next $r_i$ variables $X_1,\ldots, X_{r_i}$ processed by the sumcheck.
It claims 
\[
d_i = \sum_{\vec x\in H_{k_{i-1}}} h_i(x_1,\ldots, x_r).
\]
The verifier sends a random $\lambda_i\sample F$, which reduces  \eqref{e:zook:prev:claim} to the combined claim
\begin{align}
\label{e:zook:combined:sumcheck:claim}
\lambda_i\cdot c_{i-1} + d_i = \sum_{\vec x\in H_{k_{i-1}}} 
\left(\lambda_i \cdot G_{i-1}(\vec x) + h_i(x_1,\ldots, x_{r_i})
\right).
\end{align}

\paragraph{Step 2: Sumcheck phase.}
Both prover and verifier proceed with the $r_i$ sumcheck rounds on the combined claim \eqref{e:zook:claim:combined}.
In the course of it, the prover provides the quadratic refinement polynomials
$
\hat q_{i,1}(X), \ldots, \hat q_{i,r_i}(X), 
$ 
the verifier has sent $\vec\beta_i = (\beta_{i,1},\ldots, \beta_{i,r_i})$.
After the last sumcheck round, the claim \eqref{e:zook:combined:sumcheck:claim} is reduced to 
\begin{multline*}
\hat q_{i,r_i}(\beta_{i,r_i}) 
= \sum_{\vec x'\in H_{k_i}} \left(
\lambda_ i \cdot G_{i-1}(
\vec\beta_i,\vec x') + h_i(\vec\beta_{i})
\right)
\\
= \lambda_i\cdot \left\langle \hat p_{i-1}(\vec \beta_i, \:.\:), K^{(i-1)}(\vec\beta_i,\:.\:)\right\rangle_{H_{k_i}} + 2^{k_i - k_{i-1}}\cdot \lambda_i\cdot B_i + 2^{k_i}\cdot h_i(\vec\beta_i)
\end{multline*}
which linear in $\hat p_{i-1}(\vec\beta_{i},\:.\:)$ and the masks, now including $h_i$.

\paragraph{Step 3: Commit.}
The prover takes $\msk_{i}\in F[X]^{<q_i}$, and commits to %the masked polynomial
\[
\hat p_i = \hat p_{i-1}(\vec\beta_{i},\:.\:) + \msk_i(\:.\:)
\]
as a word from $\mathscr P_{k_{i}}^{e_i}$.
This masked commitment is required also in the last round $i=r$. 
% (For an alternative variant see Remark  \ref{rem:zook:pi:instead:pi:hat}.)
The new ensemble of polynomials, $\hat p_i, \msk_i$, and $\msk_0,\ldots, \msk_{i-1}$, $h_1,\ldots, h_i$, is subject to the linear constraint
\begin{multline}
\label{e:zook:specialized:claim}
\hat q_{i,r_i}(\beta_{i,r_i}) = 
\left\langle \hat p_i - \msk_i,  \lambda_i\cdot K^{(i-1)}(\vec\beta_i, \:.\:) \right\rangle_{H_{k_i}} 
\\
+ 2^{-r_i} \cdot \lambda_i \cdot B_{i-1} + 2^{k_i}\cdot h_i(\vec\beta_i),
\end{multline}
which are of the form \eqref{e:zook:claim}.
%except that $\msk_{i-1}$ does not appear in it.

\paragraph{Step 4: Query phase.}
In the query phase, the verifier asks the masked polynomial $\hat p_i$ for the value at the out-of-domain sample $x_{i,0}\sample F$ and uses the values of the previous $\hat p_{i-1}$ to constrain $\hat p_i - \msk_i$ at the points $x_{i,1},\ldots, x_{i,s_i}\sample D_i$.
In terms of the new ensemble, these claims are the linear constraints 
\begin{align}
\label{e:zook:out:domain:constraint}
\sum_{\vec x' \in H_{k_i}} \hat p_i(\vec x') \cdot L_{x_{i,0}}(\vec x') &= v_0,
\\
\label{e:zook:in:domain:constraint}
\sum_{\vec x'\in H_{k_i}} \left(\hat p_i(\vec x') - \msk_{i}(\vec x')\right)\cdot L_{x_{i,k}}(\vec x') &= v_{k},
\end{align}
for $k= 1,\ldots, s_i$, where $L_{x_{i_k}}$ are the univariate Lagrangians of degree $2^{k_{i-1}}$. 
These constraints are again of the form \eqref{e:zook:claim}.
%now with $\msk_{i-1}$ mentioned the first time. 


% The proximate $\hat p_1$ the multilin rep of which is
% \begin{align*}
% \langle \hat p_1,
% K_{0\lambda_1,\ldots, \lambda_{k}}
% \rangle_{H_{n-k}}
% & = q_k(\lambda_k) - \gamma\cdot h(\lambda_1,\ldots,\lambda_{k}) 
% \\
% & \quad\quad + 
% \langle \msk_1 - \msk_{0,\lambda_1,\ldots,\lambda_k},
% K_{0\lambda_1,\ldots, \lambda_{k}}\rangle_{H_{n-k}}
% \\
% \langle \hat p_1,  L_{Q_i}\rangle_{H_{n-k}} &= v_{Q_0}
% \\
% \langle \hat p_1, L_{Q_i}\rangle_{H_{n-k}} &= (\fold_{\lambda_1,\ldots, \lambda_k} \hat f_0) (Q_i).
% \end{align*}




\paragraph{Step 4: Aggregation of claims.}

The verifier samples a random $\gamma_i\sample F$ for combining the above claims \eqref{e:zook:specialized:claim}, \eqref{e:zook:out:domain:constraint} and \eqref{e:zook:in:domain:constraint} on the new ensemble into a single claim of the form \eqref{e:zook:claim}, with the following linear functionals:
\begin{align}
\label{e:zook:update:K}
K^{(i)} &=  \lambda_i\cdot K^{(i-1)}(\vec\beta_i,\:.\:) + \sum_{k=0}^{s_i} \gamma_i^{k+1} \cdot L_{x_{i,k}}(\:.\:),
\\\label{e:zook:update:Gamma}
\Gamma_{j}^{(i)}  &=  2^{r_i}\cdot\lambda_i\cdot \Lambda_j^{(i-1)}, \quad j=1,\ldots,i-1,
\\\label{e:zook:new:Gamma}
% \Gamma_i^{(i)} &= -\Big\langle \:.\:, \lambda_i\cdot K^{(i-1)}(\vec\beta_i,\:.\:) + \sum_{k=1}^{s_i} \gamma_i^{k+1}\cdot L_{x_{i,k}}(\:.\:) \Big\rangle_{H_{k_i}}   
\Gamma_i^{(i)} &= - \langle \:.\:, K^{(i)} - \gamma_i\cdot L_{x_{i,0}} \rangle_{H_{k_i}}
\\\label{e:zook:update:Lambda}
\Lambda_{j}^{(i)}  
&=  2^{r_i}\cdot\lambda_i\cdot \Lambda_j^{(i-1)}, \quad j=1,\ldots, i-1,
\\\label{e:zook:new:Lambda}
\Lambda_i^{(i)}(h_i) & = 2^{-k_i}\cdot h_i(\vec \beta_i).
\end{align}
As before, the inner product \eqref{e:zook:new:Gamma} on $\msk_i$ and the evaluation \eqref{e:zook:new:Lambda} on $h_i$ are inexpensive to compute.


\begin{rem}
\label{rem:zook:pi:instead:pi:hat}
In closer alignment with \cite{cryptoeprint:2026/391}, one may  think of the claims \eqref{e:zook:claim} in terms of the unmasked polynomial $p_i = \hat p_i - \msk_i$, rather than of the masked one, i.e.
\begin{align*}
% \label{e:zook:claim:paper}
c_{i}  &= \left\langle p_{i}, K^{(i)} \right\rangle_{H_{k_{i}}} 
+ \sum_{j=0}^{i} \Gamma_j^{(i)}(\msk_j) + \sum_{j=1}^{i} \Lambda_{j}^{(i)} (h_{j}).
\end{align*}
In Step 3, the prover commits to $\hat p_i = p_{i-1}(\vec\beta_i,\:.\:) + \msk_i(\:.\:)$ instead of $\hat p_i = \hat p_{i-1}(\vec\beta_i,\:.\:) + \msk_i(\:.\:)$, and the consistency relation for the $s_i$ in-domain queries would then be 
\begin{align*}
\sum_{\vec x'\in H_{k_i}} \left(\hat p_i(\vec x') - \msk_{i}(\vec x') + \msk_{i-1}(\vec\beta_i, \:.\:) \right)\cdot L_{x_{i,j}}(\vec x') &= v_{j}, 
\end{align*}
instead of \eqref{e:zook:in:domain:constraint}, leading to a different update rule of $\Gamma_{i-1}^{(i)}$.
\end{rem}


\subsubsection{The final opening protocol.}


% In a nutshell: 
% \begin{itemize}
% \item
% random linear combination with a random polynomial (a full mask).
% \item 
% has a separate proof of proximity, decoupled from WHIR $\longrightarrow$ multiplicative blowup by list size.
% \end{itemize}

The final claim after round $r$ is of the form
\begin{equation}
\label{e:zook:claim:r}
c_r  = \langle \hat p_r, K^{(r)} \rangle_{H_{k_r}} 
+ \sum_{j=0}^r  \Gamma_j^{(r)}(\msk_j) + \sum_{j=1}^r \Lambda_j^{(r)}(h_j).
\end{equation}
% with $K = K^{(r)}$, $\Lambda_i = \Lambda_i^{(r)}$ and $\Gamma_i = \Gamma_i^{(r)}$.
It is proven via the following elementary proof of proximity on $\hat p_r \in \mathscr P_r^{e_r}$, as a word of $C_{k_r}^{e_r}$, and the encoding of $\mathsf m = h_1\|\ldots \|h_r \| msk_0 \|\ldots \| msk_r \| msk_{zk} \in F^{m_{zk}}$ 
as word of $C_{zk}$.

\begin{protocol}
\label{prot:zook:final:opening}
The prover commits to a full entropy mask $H \sample \mathscr P_r^{e_r}$, $M\sample F^{m_{zk}}$, the latter parsed as $H_1\|\ldots\|H_r\|M_0\|\ldots\|M_r\|M_{zk}$, and claims
\begin{equation}
\label{e:zook:claim:H}
d  = \langle  H, K^{(r)} \rangle_{H_{k_r}}
+ \sum_{j=0}^r \Gamma_j^{(r)}(M_j) + \sum_{j=1}^r \Lambda_j^{(r)}(H_j).
\end{equation}
The verifier replies with $\lambda\sample F$.
The prover reveals the combined words $H' = \lambda\cdot \hat p_r + H$, $M' = \lambda\cdot \mathsf m + M$, the latter parsed as $H_1'\|\ldots\|H_r'\|M_0'\|\ldots\|M_r'\|M_{zk}'$.

% \ulrich{This needs a spec of how to commit to masks. }
The verifier queries the commitments of $\hat p_r$ and $\mathsf m$ independently at random points $x_{1}, \ldots x_{s}\sample D_r$ and $y_1,\ldots, y_s\sample D_{zk}$, respectively, and checks collinearity of the revealed values with those of $H'$ and $M'$. 
%the linear combination.
If the collinearity checks hold, and if 
\begin{equation}
\label{e:zook:claim:combined}
\lambda \cdot c_r + d  = \langle  H', K \rangle_{H_{k_r}}
+ \sum_{j=0}^r \Gamma_j(M_j') + \sum_{j=1}^r \Lambda_j(H_j'),
\end{equation}
then the verifier accepts.
\end{protocol}






\section{The complete protocol \textcolor{red}{(tbc)}}
\label{s:complete:protocol}





\section{The non-interactive STARK \textcolor{red}{(tbc)}}

\subsection{Merkle commitments \textcolor{red}{(tbc)}}

\subsection{Fiat--Shamir / BCS transform \textcolor{red}{(tbc)}}

\subsection{Example parameters \textcolor{red}{(tbc)}}

\subsection{Benchmarks \textcolor{red}{(tbc)}}


\section{Outlook v3 \textcolor{red}{(tbc)}}

% \noindent
% Provekit v3
% \begin{itemize}
% \item 
% Lean4 (feasibility study)
% \item 
% Circuits over $\mathbb Z$ and $\mathcal R = \mathbb Z[X]/(q(X))$.
% Sometimes constraints over polynomial rings really  sense. 
% Maybe some of them similarly explored in FRI-Bininus constraints?
% \item 
% WHI$\mathbb Z$: WHIR over integers.
% \item 
% Bit-integer commitments:
% Using FRI-Binius/WHIR over $\mathbb F_2$ for committing $\{0,1\}$-valued polynomials.
% Opening claim at random $\vec r\sample F_q^n$, where $q$ is the random prime,
% \[
% \sum_{\vec x\in H_n} v(\vec x)^* \cdot L_n(\vec x, \vec r)^* = a \mod q,
% \]
% via homomorphic transform of the lifted claim 
% \[
% \sum_{\vec x\in H_n} v(\vec x)^* \cdot L_n(\vec x, \vec r)^* = a^*
% \]
% to a grand product over $\mathbb F_{2^{128}}$.
% %Exploring mixed char proofs: $\mathbb Z$ combined with $\mathbb F_2$

% \item 
% Zero-knowledge for WHI$\mathbb Z$
% \end{itemize}

% \noindent
% Interesting questions around IPRS codes
% \begin{itemize}
%     \item 
%     Can we somehow lift any correlated agreement theorem from the field $\mathbb F_p$ of the construction prime to the integers $\mathbb Z$? 
%     \url{https://chatgpt.com/share/6a508b03-8e94-83ed-96eb-d58708428317}
%     \item 
%     What is the distance of an IPRS code mod a random prime?
% \end{itemize}



\bibliographystyle{alpha}
\bibliography{SNARKs}


\appendix

\section{Correlated agreement}
\label{s:CAT}

Let $\mathbb F$ be an arbitrary finite field, 
\[
\Gamma: \mathscr P_k(\mathbb F) \times \mathbb F \longrightarrow \mathbb F, %\in F[X_1,\ldots, X_n , Z],
\]
be a polynomial (in the coordinates of an arbitrary basis) of total degree $d = \deg \Gamma = 2$.


\begin{thm}[\cite{ProximityGaps,ProximityGaps:improved}] 
\label{thm:CAT:constrained}
For $C = \mathsf{RS}[F, D, k]$ any Reed--Solomon code over an arbitrary finite field $F$, with minimum distance $\delta$, and  $\Gamma:\mathscr P_k\times F\longrightarrow F$ as above.
Then, for any $\theta \in \left[\frac{\delta}{2}, 1 -\sqrt{1-\delta}\right)$, the following holds.
If the set
\[
S = \left\{z \in F\::\:  
    \begin{matrix}
    \exists p_z(X) \in \mathscr P_k,
        d((1-z)\cdot f_0 + z\cdot f_1, p_z) \leq \theta 
        \\
        \wedge \:\Gamma(p_z, z) = 0
    \end{matrix}
\right\}
\]
is of size 
\[
|S| \geq  \ell_{GS}(\theta)\cdot \left( 
\frac{2\ell_{GS}(\theta)^4}{3}\cdot (1 - \delta) %+ \frac{\theta_m |D| + 1}{|D|} 
+ 1
\right) \cdot |D|,
\]
then there exist polynomials $q_0(X)$, $q_1(X)\in\mathscr P_k$ with $d((q_0,q_1),(f_0,f_1))\leq \theta$, and $\Gamma(q_0,0)= \Gamma(q_1,1) = 0$. 
\end{thm}

The theorem is direct consequence of line decodability, also called collinearity of proximates property in \cite{Basefold:LDR}.
Given a of proximates of the random linear combinations, one for each, then a notable fraction of it lies on a line.

\section{Proofs: IOPP soundness \textcolor{red}{(tbc)}}

\subsection{Protocol \ref{prot:main}}
Let us outline the soundness properties of the protocol.
In the following we write
\begin{equation}
\ell_1 = \ell_1(\theta_1)
\end{equation}
for the Johnson list size bound on the number of proximates from $C_1$ at distance $\theta_1$.
% For both $f_1$ and $f_2$, we define $S(f_1)$ and $S(f_2)$ to be the set of $z\in F$, which distinguish all $\theta_1$- and $\theta_2$ proximates by value. 
% Note that the exceptional set 
% \begin{equation}
% \left|F\setminus S(f_i)\right| \leq \binom{\ell_i}{2} \cdot 2^{m_i},
% \end{equation} 
% %contains at most $\ell_{\theta_i}^2 \cdot 2^{m_1}$ many points, 
% the union bound of common zeros of any pair of proximates.
% We define
% \begin{equation*}
% \mathscr L_{k}\big|_{z,v} = \left\{p(X)\in \mathscr L_k \::\: p(z) = v\right\}.
% \end{equation*}
% % \begin{equation*}
% % E_z: \mathscr L_{m_1} \cap B_\theta(f_1) \longrightarrow F, \quad p(X) \mapsto p(z)
% % \end{equation*}

\paragraph{Phase \ref{i:main:phase:1}:}
The first phase is a reduction $\mathcal R_\theta \longrightarrow \mathcal R_\text{phase 1}$ depending on the verifier challenges $\vec\xi\sample \mathbb F_q^{t}$, where
\begin{align*}
\mathcal R_\text{phase 1} 
&=
\left\{
    \begin{matrix}
        f_1 \in \left(\mathbb F_q^{D}\right)^{e},% v_1\in F 
        \\
        f_2^* \in \left(\mathbb F_q^{D}\right)^{e}%, v_2\in F 
    \end{matrix}
    \::\:
    % \hspace*{-0.25cm}
    \begin{matrix}
        \exists (p_1(X),p_2(X)) \in \mathscr P_{m_1}^2 %\times \mathscr P_{m_1}
        \\
        d((p_1,p_2), (f_1,f_2^*)) \leq \theta_1,
        \\
        \text{and}
        \\
        (p_{1}(X), \vec\xi, p_{2}(X))
        \\
        \text{satisfy the lifted R1CS} 
        \\
        \text{on input }\vec{in}
        % \text{and their coefficient vectors}
        % \\
        % (\vec c_1, \vec c_2) \text{ satisfy the lifted R1CS }
        % \\
        % \text{on inputs }\vec{in} , \vec\xi 
    \end{matrix}
\right\},
\end{align*}
where as before, the distance refers to the interleaved representations of the polynomials.
The phase is a single-step reduction, and its soundness error is as follows.

\begin{lem}
Phase \ref{i:main:phase:1} in Protocol \ref{prot:main} is a round-by-round sound oracle reduction $\mathscr R_{\theta_1}\longrightarrow R_\text{phase 1}$ with soundness error
\[
\varepsilon_\text{phase 1} \leq 
\ell_1\cdot \varepsilon_\text{R1CS}, 
\]
where $\varepsilon_\text{R1CS}$ is the soundness error of the randomized R1CS.
\end{lem}

\begin{proof}
Assume that for more than $\kappa = \ell_1(\theta )\varepsilon_\text{R1CS}\cdot |\mathbb F_q|^t$ distinct verifier challenges $\vec\xi$, there exists a potential prover answer $f_2\in \mathbb F_q^{D'}$ which extends $f_1$ to a pair $(f_1,f_2^*)$ in $\mathscr R_\text{phase 1}$.
For each such challenge $\vec\xi_i$, $i=1,\ldots, \kappa$, we select one of the $\theta_1$-proximates $(p_{1}(X), p_{2}(X))$ stated by $\mathscr R_\text{phase 1}$, and obtain a list 
%$\mathscr L$ of more than $\kappa$ many $(p_{1,i}(X),\vec\xi_i, p_{2,i}(X))$.
\[
\mathscr L = \left\{
(p_{1,i}(X),\vec\xi_i , p_{2,i}(X))) \::\: 
\begin{matrix}
d(p_{1,i}(X),f_1)\leq \theta_1,
\\\text{and}
\\
(p_{1,i}(X), \vec\xi_i, p_{2,i}(X))
\\
\text{satisfy the lifted R1CS} 
\\
\text{on input }\vec{in}
\end{matrix}
\right\}
\]
with more than $\kappa$ entries. 

By the list size bound, there at most $\ell_1$ distinct $p_{1,i}(X)$ mentioned in $\mathscr L$.
One of them, $p_1(X) = p_{1,i_0}(X)$, appears more than  
\[
\kappa' = \frac{\kappa}{\ell_1} = \varepsilon_\text{R1CS}\cdot |\mathbb F_q|^t
\]
many times in the list. 
This $p_1(X)$ has more than $\kappa'$ many $\vec\xi_i, p_{2,i}(X)$ satisfying answers for the lifted randomized R1CS on input $\vec{in}$. 
Since the soundness error of the lifted R1CS equals that of the original R1CS,  $p_1(X)$ is a solution of the first stage R1CS. 
This proves that $f_1\in \mathscr R_{\theta_1}$.
$\Box$
\end{proof}

\paragraph{Phase \ref{i:main:phase:2}:}
The second phase, the zero-check, reduces $\mathcal R_\text{phase 1}$ to 
$\mathcal R_\text{phase 2}$, where 
\begin{equation*}
\mathcal R_\text{phase 2} =
\left\{
    \begin{matrix}
        f_1 \in \left(\mathbb F_q^{D_1}\right)^{e_1}, %v_1\in F 
        \\
        f_2 \in \left(\mathbb F_q^{D_k}\right)^{e_k}, %v_2\in F 
        \\
        q_1(X), \ldots, q_m(X) \in F[X]^{\leq 3}
        \\
        v_A, v_B, v_C \in F
    \end{matrix}
    \::\:
    \hspace*{-10mm}
    \begin{matrix}
        \exists (p_1(X),p_2(X)) \in \mathscr L_{m_1} \times \mathscr L_{m_2}
        \\
        d(p_1, f_1) \leq \theta_1, d(p_2, f_2)\leq \theta_2
        % \\
        % p_1(z_1) = v_1, p_2(z_2) = v_2,
        \\\text{and}
        \\
        0 = q_1(0) + q_1(1),
        \\
        q_i(\alpha_i) = q_{i+1}(0) + q_{i+1}(1),
        \\
        \text{for }i = 1,\ldots, m-1,
        \\
        q_m(\alpha_m) = (v_A \cdot v_B - v_C)\cdot eq(\vec\alpha, \vec\rho)
    \end{matrix}
\right\},
\end{equation*}
depends on the verifier challenges $\vec\rho,\vec\alpha$ (and the preceding $\vec\xi$). 
The soundness error of this reduction is 
\[
\varepsilon_\text{phase 2}  \leq \ell_1\cdot \ell_2 \cdot
    \left(
    \frac{m}{|F|}
    + m\cdot \frac{3}{|F|}
    \right), 
\]
again scaled by the number of joint $\theta_1$- and $\theta_2$-proximates of $f_1$ and $f_2$.
The components correspond to the roundwise soundness errors of that phase. 
The first term is the probability of that $\vec \rho$ is a zero of the multilinear extension of $a(\vec x)\cdot b(\vec x) - c(\vec x)$, $\vec x \in H_m$, despite being non-trivial.   
The second term is for the $m$ rounds of the sumcheck. 
Each $\frac{3}{|F|}$ stands for the probability that the specialization of the sumcheck claim to the random $\alpha_i$ is satisfied by a joint proximate $(p_1(X), p_2(X))$,
%$\in \mathscr L_{m_1}\big|_{z_1,v_1}\times \mathscr L_{m_2}\big|_{z_1,v_1}$, 
although not both claims for $X_i = 0$ and $X_i= 1$ hold. 

\paragraph{Phase \ref{i:main:phase:3}:}
The third phase reduces $\mathscr R_{phase 2}$ to the final verifier relation $\mathscr R_V$ of the protocol. 
This is 
\[
\mathscr R_{V} =
\left\{
    \begin{matrix}
        f_1 \in \left(\mathbb F_q^{D_1}\right)^{e_1}, %v_1\in F 
        \\
        f_2 \in \left(\mathbb F_q^{D_k}\right)^{e_k}, %v_2\in F 
        \\
        q_1(X), \ldots, q_m(X) \in F[X]^{\leq 3}
        \\
        v_A, v_B, v_C \in F
        \\
        g_2 \in \left(\mathbb F_q^{D_{2}}\right)^{e_{2}}, u_2\in F 
        \\
        \vdots
        \\
        g_{r-1} \in \left(\mathbb F_q^{D_{r}}\right)^{e_{r}}, u_{r-1}\in F
        \\
        g_r \in \mathscr L_{k_r} 
    \end{matrix}
    \::\:
    % \hspace*{-5mm}
    \begin{matrix}
        % \exists (p_1(X),p_2(X)) \in \mathscr L_{m_1}\big|_{z_1,v_1}\times \mathscr L_{m_2}\big|_{z_1,v_1}
        % \\
        % d(p_1, f_1) \leq \theta_1, d(p_2, f_2)\leq \theta_2
        % % \\
        % % p_1(z_1) = v_1, p_2(z_2) = v_2,
        % \\\text{and}
        \\
        0 = q_1(0) + q_1(1),
        \\
        q_i(\alpha_i) = q_{i+1}(0) + q_{i+1}(1),
        \\
        \text{for }i = 1,\ldots, m-1,
        \\
        q_m(\alpha_m) = (v_A \cdot v_B - v_C)\cdot eq(\vec\alpha, \vec\rho)
        \\\text{and}
        \\
            g_r \text{ satisfies the aggregated WHIR claim}
        \\
        \bar v = \langle K_r, g_r\rangle_{H_{k_r}}, 
        \\
        \text{which involves the specializations}
        \\
        M_i(\vec\alpha, \vec\beta \| \: .\:), 
        \\ 
        M=A,B,C, \: i=1,2
    \end{matrix}
\right\},
\]
where $g_2$, \ldots, $g_{r-1}$, are the folding proposals committed by in the first $r-1$ rounds of WHIR, together with their values $u_2$, \ldots, $u_{r-1}$ at the verifier DEEP queries $z_2,\ldots, z_{r-1}$, and $g_r$ is the final folding proposal revealed in plain. 

The soundness error of this phase is that of joint WHIR, 
\[
\varepsilon_\text{phase 3} = \varepsilon_{WHIR}(\theta_1,\ldots, \theta_r)
\]
and will be described in the following section.



\subsection{Soundness ZO0K. \textcolor{red}{(tbc)}}

For now, only few notes:

For simplicity, we assume the masks as given. 
(The full story is a bit more complicated as it needs to take into account the proximates of the masks.)
Suppose that with notable probability in $\vec\beta_i$ the prover is able to provide $\hat f_i$ which close to a poly $\hat p_i$, so that
\begin{itemize}
    \item 
    the new ensemble with $\hat p_i$ and $\msk_i$ satisfies the $\vec\beta_i$-specialized claim \eqref{e:zook:specialized:claim},
    \item 
    the folding $\mathsf{fold}_{\vec\beta_i}\hat f_{i-1}$ agrees on a density $\alpha$ subset of $D_{i-1}$ with  $\hat p_i - \msk_i$.
    (Which is caused by \eqref{e:zook:out:domain:constraint} and \eqref{e:zook:in:domain:constraint} satisfied with notable prob.)
\end{itemize}

Then, with the help of line decodability, all intermediate foldings are proximate to polys, the mv rep of which satisfy the specialized claim, up to that $\hat f_{i-1}$ has a close poly $\hat p_{i-1}$, satisfying the combined claim \eqref{e:zook:claim:combined}.



\subsection{The complete protocol \textcolor{red}{(tbc)}}

% \section{IOPP soundness}

% \subsection{Joint WHIR}

% \subsection{The complete protocol}


\section{Proofs: IOPP zero-knowledge \textcolor{red}{(tbc)}}


\subsection{ZO0K \textcolor{red}{(tbc)}}

As elaborated in Section \ref{s:zk:zerocheck}, the sumcheck polynomials of each round will be uniformly distributed over the relation of consistent polynomials, independent of the witness. 

Further, by construction the queried values (in- and out of domain) are identically distributed field elements ($F$ or $\mathbb F_q$ if the first round oracles), independent of the witness. 

The simulator can be constructed similar to the sumcheck simulator.


\subsection{The complete protocol \textcolor{red}{(tbc)}}





\section{Alternatives to ZO0K}

\subsubsection{Proposal 1: VEIL}

In its most simple variant, the followings steps need to be taken.

\subsubsection{Securing the witness queries.}
We need to randomize the witnesses $w_1$ and $w_2$ for the number of queries from $D$. (No R1CS modification necessary). 
Recall we commit the witnesses in interleaved manner, using the FFT decomposition, and thus we need to randomize the witnesses properly, so that each component is affected equally:
Thinking of $\vec w_1$ as the not yet decomposed vector, add a \textit{contiguous} interval of $e \cdot \text{num}_\text{queries}$  random values, where $e$ is the interleaving depth, at any place in the not used padding space of it. 
For $w_2$ likewise, but we may take into account that less than $e$ many components are non-zero.


\subsubsection{Isolator polynomials.}

In its most simple variant, the prover samples addionally to two random polynomials over the extension field $F$,
\[
H_1 \sample F[X_1, \ldots, X_{m_1}], \quad H_2 \sample F[X_1,\ldots, X_{m_2}],
\]
and commits them together with $w_1$ and $w_2$, respectively. 
These `isolator polynomials' are combined together with $w_1$ and $w_2$ in an additional (virtual) batching step before entering WHIR:
\begin{enumerate}
    \item 
    After we combined \eqref{e:phase3:claim:vM} and \eqref{e:phase3:claim:inputs}  into a single inner product claim $\langle K, w\rangle_{H_{m_1 + 1}}$ on the combined witness poly, the prover reveals 
    \[
    v = \langle K, H \rangle_{H_{m_1 + 1}}
    \]
    on the combined isolator polynomial $H = (1 - X_0 ) \cdot H_1 + X_0 \cdot H_2^*$.
    with $w_1$ and $w_2^*$ replaced by $H_1$ and $H_2^*$.  
    \item 
    The verifier draws a further randomness $\gamma_2$ to combine the two claims into a claim on the linear combination $w + \gamma \cdot H$. 
    This combination is not committed, the values of it are computed from those of $w$ and $H$. 
    \item 
    Run WHIR on the virtual linear combination $w + \gamma \cdot H$.
\end{enumerate}

The isolator polynomial entirely decouples the distribution of the future run of WHIR from the past (and thus the witnesses). 
WHIR is run entirely in non-zk mode, no changes here.

While the prove size increases only marginally (by one extension field element per $w_1$ query, likewise for $w_2$), prover cost does, since the prover has to hash the isolator polys. 
Moreover, we most likely hit memory bounds in this variant of VEIL.

\ulrich{%
The witness polys $w_1$ is about $1/4$-th the size as in the BN254 case, and $w_2$ is about $3/4$-th of the size.
However, $H_1$ and $H_2$ together are $3/4$-th of the entire BN254 witnesses, which is why I believe we will hit memory bounds for the largest circuits we want to prove.
}


\paragraph{A more prover friendly variant.}

To mitigate the memory bounds, we insert isolator polynomials in the first round of WHIR instead.
\ulrich{This is a hacky solution, which requires additional sumcheck blinders, I am not a fan of it.}

\begin{enumerate}
\item
The prover only samples a single
\[
H \sample C_1
\]
from the base code of $w_1$ and $w_2^*$, and commits either together with $w_1$. 
This is much less overhead (both compute and memory). 

\item 
Modified first round of WHIR:
\begin{enumerate}
\item
The sumcheck phase is randomized in the usual manner, using 
\[
h_1(X),\ldots, h_{k_1}(X)\sample F[X]^{\leq 3}.
\]
These are committed separately, and in zk manner. 

\item
After specializing the sumcheck claim to $X_1 = \beta_1,\ldots, X_k =\beta_{k_1}$, the prover reveals 
\[
v = \langle K_0(\beta_1,\ldots, \beta_{k_1}, \:.\:), H\rangle_{H_{m_1 - k_1}},
\]
and the claims for $w(\beta_1,\ldots, \beta_{k_1},\:.\:)$ and $H$ are combined via another verifier randomness $\gamma_2\sample F$.
\item 
We now proceed with the virtual linear combination $w(\beta_1,\ldots, \beta_{k_1}, \:.\:) + \gamma_2\cdot H$.
(The query phase, the aggregation phase.)
\end{enumerate}
\end{enumerate}

The value $h_1(\beta_1) +\ldots + h_{k_1}(\beta_{k_1})$ is proven separately. 



\end{document}
